%%%%%%%% Manuscript v5.2 (anonymous review copy) %%%%%%%%
% TARGET_YEAR      = 2027  (ICML 2027; icml.cc/Conferences/2027, its CallForPapers page and the icml2027.zip style URL returned 404 on 2026-09-27)
% TEMPLATE_YEAR    = 2026  (official ICML 2026 style kit, review/anonymous mode, unmodified)
% SUBMISSION_READY = false
% Build: scripts/build_paper.sh (copies the official kit, sha256-checked, into paper/build/;
%        the style files are not vendored in the repository).
% All numbers in the text come from paper/generated/numbers.tex; all tables and the figure
% from paper/generated/*.tex, produced by scripts/make_paper_assets.py from results/raw.
%%%%%%%%%%%%%%%%%%%%%%%%%%%%%%%%%%%%%%%%%%%%%%%%%%%%%%%%%%%%%%%%%%%%%%%%%%%%%%%%%%%%%%

\documentclass{article}

\usepackage{microtype}
\usepackage{graphicx}
\usepackage{subcaption}
\usepackage{booktabs}
\usepackage{placeins}
\usepackage{hyperref}

\newcommand{\theHalgorithm}{\arabic{algorithm}}

% Initial blind version submitted for review (official ICML 2026 style; see header):
\usepackage{icml2026}

\usepackage{amsmath}
\usepackage{amssymb}
\usepackage{mathtools}
\usepackage{amsthm}
\usepackage[capitalize,noabbrev]{cleveref}
\usepackage{pgfplots}
\usepgfplotslibrary{groupplots}
\pgfplotsset{compat=1.17}

\theoremstyle{plain}
\newtheorem{theorem}{Theorem}[section]
\newtheorem{proposition}[theorem]{Proposition}
\newtheorem{lemma}[theorem]{Lemma}
\theoremstyle{definition}
\newtheorem{definition}[theorem]{Definition}
\newtheorem{assumption}[theorem]{Assumption}
\theoremstyle{remark}
\newtheorem{remark}[theorem]{Remark}

% Visible marker for evidence that does not exist yet (never a placeholder number).
\newcommand{\todo}[1]{\textcolor{red}{\textbf{[TODO: #1]}}}
\newcommand{\notrun}{\textsc{not run}}

% AUTO-GENERATED by scripts/make_paper_assets.py from results/raw -- do not edit by hand.
\newcommand{\AdapterParams}{\ensuremath{23{,}397}}
\newcommand{\AdapterStep}{\ensuremath{0.031}}
\newcommand{\CompErrEight}{\ensuremath{0.026}}
\newcommand{\CompErrOne}{\ensuremath{0.19}}
\newcommand{\CompInv}{\ensuremath{1.7\times10^{-15}}}
\newcommand{\CompRK}{\ensuremath{6.1\times10^{-15}}}
\newcommand{\CompSeconds}{\ensuremath{0.11}}
\newcommand{\CompSlope}{\ensuremath{-0.95}}
\newcommand{\CompSlopeEq}{\ensuremath{-0.93}}
\newcommand{\ConvCodeRight}{\ensuremath{6.5\times10^{-19}}}
\newcommand{\ConvShiftIrregular}{\ensuremath{3.7\times10^{-3}}}
\newcommand{\ConvShiftUniform}{\ensuremath{8.7\times10^{-19}}}
\newcommand{\DecABMax}{\ensuremath{1.26}}
\newcommand{\DecABMin}{\ensuremath{1.22}}
\newcommand{\DecDeltaAbsMax}{\ensuremath{0.62}}
\newcommand{\DecDeltaAbsMin}{\ensuremath{0.60}}
\newcommand{\DecDeltaMax}{\ensuremath{+0.50}}
\newcommand{\DecDeltaMin}{\ensuremath{+0.46}}
\newcommand{\DecIdentity}{\ensuremath{4.6\times10^{-15}}}
\newcommand{\DecMMax}{\ensuremath{+0.011}}
\newcommand{\DecMMin}{\ensuremath{-0.016}}
\newcommand{\DecMabsMax}{\ensuremath{0.028}}
\newcommand{\DecMabsMin}{\ensuremath{0.023}}
\newcommand{\DecMendMax}{\ensuremath{0.0049}}
\newcommand{\DecMendMin}{\ensuremath{0.0005}}
\newcommand{\DecMismatch}{\ensuremath{0}}
\newcommand{\DecQabsMax}{\ensuremath{0.023}}
\newcommand{\DecQabsMin}{\ensuremath{0.013}}
\newcommand{\DecRepErrMax}{\ensuremath{0.59}}
\newcommand{\DecRepErrMin}{\ensuremath{0.56}}
\newcommand{\DecSignMin}{\ensuremath{98}}
\newcommand{\DecWMax}{\ensuremath{+0.49}}
\newcommand{\DecWMin}{\ensuremath{+0.48}}
\newcommand{\DecWabsMax}{\ensuremath{0.62}}
\newcommand{\DecWabsMin}{\ensuremath{0.60}}
\newcommand{\DecWall}{\ensuremath{7.7}}
\newcommand{\EoneBilinSlope}{\ensuremath{-2.02}}
\newcommand{\EoneComplexBilin}{\ensuremath{0.24}}
\newcommand{\EoneEulerMeight}{\ensuremath{0.025}}
\newcommand{\EoneEulerMone}{\ensuremath{0.21}}
\newcommand{\EoneEulerSlope}{\ensuremath{-1.03}}
\newcommand{\EoneNodtMeight}{\ensuremath{1.3}}
\newcommand{\EoneNodtMtwo}{\ensuremath{0.48}}
\newcommand{\EoneRKmax}{\ensuremath{1.6\times10^{-13}}}
\newcommand{\EoneRealBilin}{\ensuremath{9.2\times10^{-3}}}
\newcommand{\EoneZohMax}{\ensuremath{6.3\times10^{-15}}}
\newcommand{\EthreeExactChange}{\ensuremath{1.5\times10^{-16}}}
\newcommand{\EthreeExactErrEight}{\ensuremath{0.032}}
\newcommand{\EthreeMeanChange}{\ensuremath{0.30}}
\newcommand{\EthreeMeanErrEight}{\ensuremath{0.28}}
\newcommand{\EthreeMeanErrOne}{\ensuremath{0.060}}
\newcommand{\EthreeRMEneg}{\ensuremath{21}}
\newcommand{\EthreeRMEslope}{\ensuremath{-0.076}}
\newcommand{\EthreeRiemannChange}{\ensuremath{0.026}}
\newcommand{\EthreeSumChange}{\ensuremath{1.5}}
\newcommand{\EthreeTrapzChange}{\ensuremath{0.011}}
\newcommand{\EtwoCrossMax}{\ensuremath{+0.42}}
\newcommand{\EtwoCrossMin}{\ensuremath{-1.37}}
\newcommand{\EtwoEulerArtFrac}{\ensuremath{0.53}}
\newcommand{\EtwoEulerCrossFrac}{\ensuremath{+0.27}}
\newcommand{\EtwoHoldErrEight}{\ensuremath{0.030}}
\newcommand{\EtwoHoldErrOne}{\ensuremath{0.23}}
\newcommand{\EtwoNodtArtFrac}{\ensuremath{1.22}}
\newcommand{\EtwoNodtCrossFrac}{\ensuremath{-0.25}}
\newcommand{\EtwoSlope}{\ensuremath{-0.98}}
\newcommand{\EtwoZohArt}{\ensuremath{1.7\times10^{-13}}}
\newcommand{\FacAll}{\ensuremath{1.5\times10^{-16}}}
\newcommand{\FacDt}{\ensuremath{0.30}}
\newcommand{\FacDtTime}{\ensuremath{0.049}}
\newcommand{\FacDtZoh}{\ensuremath{0.30}}
\newcommand{\FacDtZohTime}{\ensuremath{0.011}}
\newcommand{\FacNone}{\ensuremath{0.46}}
\newcommand{\FacTime}{\ensuremath{0.25}}
\newcommand{\FacZoh}{\ensuremath{0.48}}
\newcommand{\FirstRunSeconds}{\ensuremath{20.6}}
\newcommand{\HfourEulerFrac}{\ensuremath{31}}
\newcommand{\HsevenOrigSlope}{\ensuremath{+0.42}}
\newcommand{\NfourFloor}{\ensuremath{1.1\times10^{-4}}}
\newcommand{\NoneNaive}{\ensuremath{8.0\times10^{-4}}}
\newcommand{\NthreeEulerLarge}{\ensuremath{4.6}}
\newcommand{\NthreeEulerSmall}{\ensuremath{3.5\times10^{-3}}}
\newcommand{\NtwoAbar}{\ensuremath{-0.977}}
\newcommand{\NtwoBilArt}{\ensuremath{1.1}}
\newcommand{\NtwoEulerArt}{\ensuremath{170}}
\newcommand{\NunitsPrimary}{\ensuremath{32}}
\newcommand{\RealDevMSE}{\ensuremath{0.012}}
\newcommand{\RealDevMax}{\ensuremath{0.84}}
\newcommand{\RealDiscSeight}{\ensuremath{0.011}}
\newcommand{\RealDiscSeightCIhi}{\ensuremath{0.011}}
\newcommand{\RealDiscSeightCIlo}{\ensuremath{0.010}}
\newcommand{\RealDiscStwo}{\ensuremath{6.1\times10^{-3}}}
\newcommand{\RealEvalSec}{\ensuremath{24}}
\newcommand{\RealJoneMSE}{\ensuremath{1.1}}
\newcommand{\RealLabelSD}{\ensuremath{2.35}}
\newcommand{\RealLayerOneDp}{\ensuremath{1.8\times10^{-14}}}
\newcommand{\RealLayerOneFp}{\ensuremath{7.2\times10^{-6}}}
\newcommand{\RealLayerTwoMax}{\ensuremath{7.2\times10^{-3}}}
\newcommand{\RealLayerTwoMin}{\ensuremath{4.1\times10^{-3}}}
\newcommand{\RealMSEczero}{\ensuremath{0.0167}}
\newcommand{\RealMSEheight}{\ensuremath{0.0156}}
\newcommand{\RealMSEseight}{\ensuremath{0.0160}}
\newcommand{\RealPoolCzero}{\ensuremath{0.063}}
\newcommand{\RealPoolLabelAbs}{\ensuremath{1.59}}
\newcommand{\RealPoolMeanHeight}{\ensuremath{0.56}}
\newcommand{\RealPoolResHeight}{\ensuremath{0.063}}
\newcommand{\RealPoolTimeHeight}{\ensuremath{0.068}}
\newcommand{\RealRelDiffHeight}{\ensuremath{-6.5}}
\newcommand{\RealRelDiffSeight}{\ensuremath{-4.1}}
\newcommand{\RealRelDiffStwo}{\ensuremath{-3.0}}
\newcommand{\RealRunNatCzero}{\ensuremath{0.053}}
\newcommand{\RealRunNatSeight}{\ensuremath{1.07}}
\newcommand{\RealRunResSeight}{\ensuremath{0.054}}
\newcommand{\RealStep}{\ensuremath{1500}}
\newcommand{\RealTrainSec}{\ensuremath{97}}
\newcommand{\RtwoCzeroRatioMax}{\ensuremath{1.04}}
\newcommand{\RtwoCzeroRatioMin}{\ensuremath{0.98}}
\newcommand{\RtwoDataWall}{\ensuremath{40}}
\newcommand{\RtwoDiscTrajMin}{\ensuremath{1.6\times10^{-3}}}
\newcommand{\RtwoDlomin}{\ensuremath{0.512}}
\newcommand{\RtwoDmax}{\ensuremath{0.602}}
\newcommand{\RtwoDmin}{\ensuremath{0.578}}
\newcommand{\RtwoEqEqual}{\ensuremath{\text{exactly equal}}}
\newcommand{\RtwoEqMax}{\ensuremath{0}}
\newcommand{\RtwoEvalWall}{\ensuremath{143}}
\newcommand{\RtwoFpEstMax}{\ensuremath{5.2\times10^{-6}}}
\newcommand{\RtwoFpRelMax}{\ensuremath{3.8\times10^{-4}}}
\newcommand{\RtwoFpTrajMax}{\ensuremath{1.8\times10^{-5}}}
\newcommand{\RtwoFratioMax}{\ensuremath{14.1}}
\newcommand{\RtwoFratioMin}{\ensuremath{11.5}}
\newcommand{\RtwoInvMaxPct}{\ensuremath{9.5}}
\newcommand{\RtwoInvN}{\ensuremath{4}}
\newcommand{\RtwoInvTot}{\ensuremath{60}}
\newcommand{\RtwoPoolMeanMax}{\ensuremath{0.59}}
\newcommand{\RtwoPoolMeanMin}{\ensuremath{0.57}}
\newcommand{\RtwoPoolResMax}{\ensuremath{0.096}}
\newcommand{\RtwoPoolResMin}{\ensuremath{0.066}}
\newcommand{\RtwoPoolTimeMax}{\ensuremath{0.094}}
\newcommand{\RtwoPoolTimeMin}{\ensuremath{0.068}}
\newcommand{\RtwoRHeightMax}{\ensuremath{+6.7}}
\newcommand{\RtwoRHeightMin}{\ensuremath{-4.9}}
\newcommand{\RtwoRJoneDev}{\ensuremath{-20.9}}
\newcommand{\RtwoRJoneMax}{\ensuremath{-21.2}}
\newcommand{\RtwoRJoneMin}{\ensuremath{-22.6}}
\newcommand{\RtwoRJonePilotOld}{\ensuremath{-6.4}}
\newcommand{\RtwoRSeightMax}{\ensuremath{+3.1}}
\newcommand{\RtwoRSeightMin}{\ensuremath{-6.2}}
\newcommand{\RtwoRatioMax}{\ensuremath{11.5}}
\newcommand{\RtwoRatioMin}{\ensuremath{9.7}}
\newcommand{\RtwoResOverMax}{\ensuremath{0.019}}
\newcommand{\RtwoResOverMin}{\ensuremath{0.007}}
\newcommand{\RtwoSeightFold}{\ensuremath{5.1}}
\newcommand{\RtwoSeightMax}{\ensuremath{0.022}}
\newcommand{\RtwoSeightMin}{\ensuremath{4.3\times10^{-3}}}
\newcommand{\RtwoStepsMax}{\ensuremath{1800}}
\newcommand{\RtwoStepsMin}{\ensuremath{800}}
\newcommand{\RtwoTestWall}{\ensuremath{16}}
\newcommand{\RtwoTrainCPU}{\ensuremath{364}}
\newcommand{\RtwoTrainWall}{\ensuremath{188}}


\icmltitlerunning{Where Sampling-Grid Dependence Enters Selective SSMs}

\begin{document}

\twocolumn[
  \icmltitle{Where Sampling-Grid Dependence Enters Selective State-Space Models}

  \begin{icmlauthorlist}
    \icmlauthor{Anonymous Authors}{anon}
  \end{icmlauthorlist}

  \icmlaffiliation{anon}{Anonymous Institution, Anonymous City, Anonymous Country}

  \icmlcorrespondingauthor{Anonymous Authors}{anonymous@example.com}

  \icmlkeywords{state-space models, selective SSM, discretization, irregular sampling, zero-order hold}

  \vskip 0.3in
]

\printAffiliationsAndNotice{}

\begin{abstract}
A sequence model can change its output when only the time grid of the same physical input changes; we ask where this dependence enters a selective state-space model (SSM), separating interval splitting, which keeps the held input path, from new observations, which change it.
For one layer, a physical step, exact zero-order hold and coefficients that see only the held value give split invariance; in a stack, later layers receive an internal signal that the refined grid samples differently.
On one selective SSM with a physical step (official TIDES code) and four seeds on one synthetic task, we separate a model-coupling change at common times, replicated in every seed; a token-mean readout change under a density change, mostly the weighted-average identity of the readout; and a native-versus-resampled task-risk contrast, which changes sign across seeds.
One checkpoint on UCI HAR (activity recognition, nine subjects) shows the same separation: under a pooled splitting change of $\HarCtRelAll$ in class-centered logits, the change at the common time points is $\HarCtRelEnd$, an order of magnitude below the readout term over the added tokens; the two are not additive shares.
The pre-registered cross-entropy difference between token-mean and time-weighted pooling is $\HarD$ nats with a subject-level interval that includes zero, so no risk difference is established; two further seeds on the same subjects give $\RviiiDA$ and $\RviiiDB$ nats as an unpooled sensitivity check. We propose no new architecture.
\end{abstract}

%=====================================================================================
\section{Introduction}
\label{sec:intro}

Continuous-time state-space layers are applied to sequences by discretizing a linear ODE over each interval between observations.
When the step is the physical elapsed time, the recurrence inherits the semantics of the continuous system, which lets S4, LSSL and S5 change the sampling rate at test time or consume irregular timestamps~\citep{gu2021s4,gu2021lssl,smith2022s5}.
Selective SSMs such as Mamba compute the step from the current token~\citep{gu2023mamba,dao2024mamba2}, so it acts as a content gate; TIDES restores a physical step and moves selectivity to the state matrix~\citep{soydan2026tides}.

We do not propose another architecture.
We ask a measurement question: \emph{when only the time grid of the same physical input path changes, where does the output change enter?}
\emph{Splitting} a held interval into $m$ sub-steps that repeat the held value leaves the held input path unchanged; we call an output change at the common times an \emph{artifact}, in the operational sense of a violation of this pre-declared invariance.
\emph{Adding observations} of the same signal changes the observed values and hence the reconstructed input path, so part of the change is a response to a different input; we make no claim about task-relevant or Bayesian information.

\paragraph{Known and not claimed.}
The physical $\Delta t$, exact zero-order-hold (ZOH) formulas and flow composition are standard~\citep{gu2021lssl,smith2022s5}; that the released Mamba code uses $\bar B=\Delta B$ (``exponential-Euler'') is documented by \citet{lahoti2026mamba3}; that a continuous-time reading of selective SSMs needs $\Delta\propto dt$ is shown by \citet{cirone2024theoretical}; keeping $\Delta$ physical with selectivity on the decay rate is the design of \citet{soydan2026tides}; a sample mean weighting samples rather than time is elementary.

\paragraph{Contributions.}
None is an architecture.
\begin{enumerate}
  \item A protocol that separates splitting from new observations, with an exact path/artifact decomposition and explicit denominators (\cref{sec:setup}).
  \item Two different causes, stated with their assumptions: resampling of internal signals between layers, which changes a stack of exact-update layers under splitting although the held input is unchanged, with an $O(\text{mesh})$ upper bound for a scalar linear cascade (\cref{sec:coupling}), and sample-count readouts, which weight tokens instead of time (\cref{sec:readout-analysis}); toy checks in \cref{sec:toy}.
  \item A fixed-protocol seed repetition on the official TIDES code that separates a model-coupling effect, a pooled-weighting effect and a task-risk effect, with cost; an exact signed decomposition shows that the token-mean readout effect is almost entirely a weighted-average identity (\cref{sec:real-syn}); and one pre-registered evaluation of a checkpoint trained on real sensor windows, with a fixed-checkpoint diagnostic that separates the common-time output change from the readout's averaging over added time points, and a two-seed sensitivity check (\cref{sec:har}).
  \item A check that the TIDES paper equations and the released code pair inputs with different intervals (\cref{sec:convention}).
\end{enumerate}

%=====================================================================================
\section{Related Work}
\label{sec:related}

\paragraph{Discretization in structured SSMs.}
S4 uses the bilinear transform and reports a test-time resolution change obtained by changing its internal step~\citep{gu2021s4}; LSSL handles irregular data by discretizing with a different timescale~\citep{gu2021lssl}; S4D gives bilinear and ZOH rules~\citep{gu2022s4d}.
S5 uses ZOH, $\bar\Lambda=e^{\Lambda\Delta}$, $\bar B=\Lambda^{-1}(\bar\Lambda-I)\tilde B$, rescales $\Delta$ globally for zero-shot 8\,kHz evaluation of a 16\,kHz model, and supplies a per-step $\Delta_t$ for an irregular pendulum task, with ablations that drop the interval or append it as a feature~\citep{smith2022s5}.
S4ND notes that only the relative rate matters~\citep{nguyen2022s4nd}; event-camera SSMs evaluate at different frequencies~\citep{zubic2024event}.
These are LTI layers.

\paragraph{Selective SSMs.}
Mamba makes $\Delta,B,C$ input-dependent and states the ZOH rule~\citep{gu2023mamba}; Mamba-2 keeps its discretization~\citep{dao2024mamba2}.
Mamba-3 derives right-endpoint ZOH for time-varying systems, names the released rule exponential-Euler, and proposes an exponential-trapezoidal rule that is second order only if its data-dependent weight is $\tfrac12+O(\Delta_t)$~\citep{lahoti2026mamba3}.
Liquid-S4 makes the transition input-dependent~\citep{hasani2022liquid}.

\paragraph{Physical time in selective SSMs.}
TIDES keeps $\tilde\Delta\equiv\Delta$, makes $\mathrm{Re}\,\Lambda$, $B$ and $C$ input-dependent, and uses $\bar B_k=\Lambda_k^{-1}(\bar\Lambda_k-I)B_k$~\citep{soydan2026tides}.
Its Fading Flash diagnostic rescales a global step for fixed-length sequences and shows that a Mamba surrogate with $\Delta$ as an input channel fails outside the training range; its random-drop study removes observations and reports accuracy.
In the sections we read it neither refines a fixed held path nor decomposes output changes; it is the closest prior work, and we evaluate its official code, not a re-implementation, treating its architecture as prior art.

\paragraph{Continuous-time models.}
Neural CDEs are driven by an interpolation of the observations~\citep{kidger2020ncde}; Log-NCDEs build on NCDEs, which are robust to irregular sampling~\citep{walker2024logncde}; \citet{cirone2024theoretical} read selective SSMs as linear CDEs in a limit where $\delta$ plays the role of $dt$; oscillatory SSMs study implicit and IMEX integration~\citep{rusch2024linoss}.
We did not find, in these sources, a refinement of the same held path for selective SSMs that separates splitting from new observations and treats readouts and inter-layer resampling as entry points; this is absence of evidence, not proof of novelty.

%=====================================================================================
\section{Problem Setup}
\label{sec:setup}

\paragraph{Model and holds.}
The toy is a scalar-input selective model with $N$ diagonal modes,
\begin{equation}
\dot h_n=g(u)\big(\lambda_n h_n+B_n(u)\,u\big),\qquad y=\mathrm{Re}\textstyle\sum_nC_n(u)h_n,
\label{eq:ct}
\end{equation}
with $\mathrm{Re}\,\lambda_n<0$, a gate $g(u)=\mathrm{softplus}(wu+\beta)$ shared across modes, and affine $B(u),C(u)$.
A grid $0=t_0<\dots<t_K=T$ has steps $dt_k$ and mesh $\delta$.
We use the right-endpoint hold $\tilde u(t)=u(t_k)$ on $(t_{k-1},t_k]$, which is the convention of the Mamba recurrence and of the TIDES code (\cref{sec:convention}).
A discrete variant is $h_k=\bar A_kh_{k-1}+\bar B_ku_k$ with $\Delta_k=dt_k\,g(u_k)$ (``$\Delta{=}dt\,g$'') or $\Delta_k=\tau g(u_k)$ with a fixed $\tau$ (``$\Delta{=}\tau g$''), and exact ZOH, Euler-$B$ ($\bar B=\Delta B$) or bilinear rules; $\tau$ is the base step.

\paragraph{Grid changes and decomposition.}
\emph{Splitting} divides base intervals into $m$ equal sub-steps that all receive the held value; \emph{new observations} give each new sub-time the value of the continuous signal.
With $D_m$ the outputs of a variant on grid $m$ and $Y_m$ the continuous-time solution of \eqref{eq:ct} on the held path of grid $m$, the change $T_m=D_m-D_1$ at the common base times splits exactly as
\begin{equation}
T_m=\underbrace{(Y_m-Y_1)}_{P_m\text{: path}}+\underbrace{\big[(D_m-Y_m)-(D_1-Y_1)\big]}_{R_m\text{: artifact}},
\label{eq:decomp}
\end{equation}
so $\|T_m\|^2=\|P_m\|^2+\|R_m\|^2+2\langle P_m,R_m\rangle$.
Under splitting $P_m=0$.
Because the cross term can be large and of either sign, we never report a ratio of norms as a share or as a causal contribution.

\paragraph{Readouts.}
We distinguish sample-count readouts ($\sum_ky_k$, $\frac1K\sum_ky_k$), time-weighted endpoint sums (right-endpoint $\frac1T\sum_ky_k\,dt_k$ and trapezoidal), and the time average $\frac1T\int_0^Ty\,dt$ of the output over the held path, which the toy computes in closed form for exact-ZOH states (the ``exact integral'').

%=====================================================================================
\section{Where Grid Dependence Enters: Analysis}
\label{sec:analysis}

\subsection{Single layer}
\label{sec:single}

\begin{proposition}[standard]
\label{prop:split}
Fix a base interval $(t_{k-1},t_k]$ and suppose
(i) every sub-step receives the held value $u_k$;
(ii) every coefficient ($g,B,C,\lambda$) is a function of the held value only;
(iii) the step is (sub-step length)$\,\cdot g(u_k)$; and
(iv) $\bar A,\bar B$ are the exact ZOH formulas.
Then, from the same state at $t_{k-1}$, the state and output at $t_k$ are the same for every $m$; by induction the outputs at all base times are split-invariant.
\end{proposition}
\begin{proof}
Each sub-step is the exact flow $\Phi_s$ of $\dot h=g(u_k)(\lambda h+B(u_k)u_k)$ over its length $s$, and $\Phi_{s_1}\circ\Phi_{s_2}=\Phi_{s_1+s_2}$.
\end{proof}

The proposition is textbook; its use is that each entry point violates one assumption.
$\Delta{=}\tau g$ violates (iii): $m$ sub-steps apply $m$ times the decay.
Euler-$B$ violates (iv), with local error $|\Delta Bu|\,|\varphi_1(\Delta\lambda)-1|\le|Bu|\Delta^2|\lambda|/2$, $\varphi_1(z)=(e^z-1)/z$, since $e^z-1-z=z^2\int_0^1(1-s)e^{sz}ds$ gives $|\varphi_1(z)-1|\le|z|/2$ for $\mathrm{Re}\,z\le0$.
(ii) fails whenever a coefficient sees anything besides the held value: a token-indexed convolution (Mamba's \texttt{d\_conv}${=}4$), time or position features, state-dependent transitions, or training-mode normalization statistics over time.
Bilinear is stable for $\mathrm{Re}\,z<0$, but for real $z<-2$ its transition is negative (sign-alternating transients) and tends to $-1$ instead of $e^z\to0$ (inaccurate stiff decay); its held-input fixed point is exact.

\subsection{Several layers: internal resampling}
\label{sec:coupling}

In a stack, layer~2 is driven by the continuous output of layer~1, but a layerwise recurrence gives it only endpoint samples held over each step.
Splitting the input grid changes which internal samples it sees, without adding any external observation.
The minimal instance is
\begin{equation}
\dot h_1=-a h_1+b\,u(t),\qquad \dot h_2=-c\,h_2+d\,h_1(t).
\label{eq:cascade}
\end{equation}

\begin{proposition}[standard]
\label{prop:cascade}
Let $u$ be held on a grid with mesh $\delta$, $|u|\le U$, $a\ge0$, $c>0$ and $h_1(0)=h_2(0)=0$.
(a) The exact coupled solution of \eqref{eq:cascade} at base times is split-invariant.
(b) The layerwise computation that is exact in layer~1 and holds $h_1(t_k)$ on $(t_{k-1},t_k]$ in layer~2 has absolute error $|e_k|\le|b\,d|\,U\,\delta\,e^{c\delta}/c$ at every grid time, for every $T$ and also for $a=c$.
\end{proposition}
Proof in \cref{app:proof}.

\begin{remark}[scope of \cref{prop:cascade}]
\label{rem:scope}
The argument is elementary and has not been checked by a third party.
A nonzero $h_1(0)$ adds $a|h_1(0)|$ to $2|b|U$; a nonzero $h_2(0)$ does not enter.
Uniformity in $T$ needs $c>0$: as $c\to0$ the sum is bounded only by $t_k$ and the bound grows linearly in time; $a<0$ (an unstable first layer) is not covered.
The statement is for a linear time-invariant scalar cascade; it gives no guarantee for input-dependent coefficients, nonlinear maps between layers, or trained deep models, where the Lipschitz and contraction constants it uses are not controlled.
\end{remark}

Conditions (ii)--(iv) of \cref{prop:split} concern each layer's update and can hold in every layer, but condition (i) holds only for the first: the input of layer~2 is the output of layer~1, which is not constant within a base interval, so a stacked recurrence on the observation grid sees a different layer-2 input after splitting.
The resulting change is a reconstruction error of an \emph{internal} signal, not a change of the observed input.
\Cref{prop:cascade}(b) is an upper bound of order $\delta$ for the stated scalar cascade; it does not show that the error is nonzero or of first order for every input (for $u\equiv0$ it is zero), and it is not a theorem about deep selective networks.
Two quantities should be kept apart.
The \emph{invariance defect}, the change of the base-time outputs under splitting, vanishes if the inter-layer signal is represented exactly or if the stack runs on a fixed grid that does not depend on the observations (for example the training grid after resampling).
The \emph{approximation error} of the layerwise computation relative to the continuous coupled system does not vanish on a fixed grid; for the cascade it is bounded as in \cref{prop:cascade}(b) with the mesh of that grid.

\subsection{Readouts}
\label{sec:readout-analysis}
The time average of the output over the held path is additive over sub-intervals, so it is split-invariant whenever the state trajectory is (for example under the conditions of \cref{prop:split}); endpoint sums converge to it only where the grid is refined; sample-count readouts weight samples, not time, so the sample mean converges to a density-weighted average.
If the first half of $[0,T]$ is split by $m$ and the second half is not, the token mean of outputs that merely repeat the base-grid outputs differs from their base-grid mean by $\frac{m-1}{2(m+1)}(A-B)$, where $A$ and $B$ are the mean base-grid outputs over the two halves ($\frac{7}{18}(A-B)$ for $m{=}8$); this weighted-average identity involves no property of the model.

\subsection{Interval conventions in the TIDES paper and code}
\label{sec:convention}

The TIDES paper holds $u_k$ forward on $[t_k,t_{k+1})$ with the state in $y_k$ not yet containing $u_k$, whereas the pinned code holds $u_k$ backward on $(t_{k-1},t_k]$ with the state in $y_k$ containing it; both are exact ZOH solutions of different input reconstructions, and on an irregular grid the difference is not a relabeling (checked with untrained weights: \cref{app:convention}).
The released code coincides with the right-endpoint convention used throughout this paper, so our splitting protocol repeats the right-endpoint value accordingly; we make no claim that this affects TIDES' results.

%=====================================================================================
\section{Trained Official TIDES: Synthetic Seed Repetition and One Real-Sensor Evaluation}
\label{sec:real}

\subsection{Synthetic task: fixed-protocol seed repetition}
\label{sec:real-syn}

\paragraph{Setup.}
\emph{Model:} official TIDES at a pinned commit (MIT), $\AdapterParams$ parameters, two blocks, input-dependent $\mathrm{Re}\,\Lambda$, $B$, $C$, exact ZOH, causal, no convolution, a linear per-step head; CPU, float32, 2 threads.
\emph{Task:} regression of the teacher \eqref{eq:ct} with fixed parameters, driven by sine-sum inputs, over $T{=}8$ on a uniform training grid of $K{=}32$ steps; labels are the teacher output on the continuous path at the 32 grid times (label s.d.\ $\RealLabelSD$); trajectory identifiers are split into 512/128 train/calibration before any grid is built, with normalization from the training split only.
\emph{Training:} Adam, 2000 steps, batch 32; calibration MSE every 100 steps; the lowest-calibration-MSE state is saved, hashed and frozen before any test trajectory is generated.
A development run with seed 0 (\cref{app:pilot}) chose the questions below; the repetition was then fixed: seeds 1--4 for the initialization and minibatch order, train and calibration data hash-identical to the development run, and the repetition code reproduces every tensor of the development checkpoint (maximum absolute difference $\RtwoEqMax$).
All seeds share one new test set of 256 trajectories from the same generator, with identifiers disjoint from all earlier ones; no seed was replaced or added.

\paragraph{Conditions, arms and estimands.}
The held-path refinements are C0 (training grid), S2/S4/S8 (every interval split with the held value repeated) and H8 (only the first half split by 8, a density change on the same held path); J1 ($K$ random observation times) changes the observed path and is reported separately.
\emph{Native} evaluation feeds the condition's grid with physical step sizes; \emph{resampled} evaluation carries the last observation forward onto the training grid using only observed times and values.
On S and H conditions resampling reproduces the C0 input exactly, so its zero discrepancy is an identity control, not evidence of superiority.
For a sequence-level readout let $\bar y_i=T^{-1}\!\int_0^Ty_i\,dt$ be the pooled label (Simpson's rule on the teacher's RK4 nodes) and $\hat y_{i,k}(c)$ the native outputs on condition $c$ with steps $dt_k(c)$ and $L_c$ tokens.
We compare the token mean $p^{\mathrm{mean}}_i(c)=L_c^{-1}\sum_k\hat y_{i,k}(c)$, the pooling of the TIDES classifier head (here applied to the outputs of a linear per-step head, which equals pooling the features first), with the time-weighted sum $p^{\mathrm{time}}_i(c)=T^{-1}\sum_k\hat y_{i,k}(c)\,dt_k(c)$ of the \emph{same} outputs.
The latter is a pooling-only ablation, not part of the official model and not an exact integral; the resampled arm pools its $K$ outputs by $p^{\mathrm{res}}_i(c)=K^{-1}\sum_k\hat y^{\mathrm{res}}_{i,k}(c)$. On the uniform C0 grid all three equal $p_i(\mathrm{C0})$; the pooled error of an arm is $n^{-1}\sum_i|p_i(c)-\bar y_i|$.
The primary estimand, chosen after the development run, is the pooled-readout effect
$D_s=\frac1n\sum_i\big(|p^{\mathrm{mean}}_{s,i}(\mathrm{H8})-p_{s,i}(\mathrm{C0})|-|p^{\mathrm{time}}_{s,i}(\mathrm{H8})-p_{s,i}(\mathrm{C0})|\big)$ in label units.
Secondary estimands are the final-output discrepancy under S8, $\|\hat y(\mathrm{S8})-\hat y(\mathrm{C0})\|_2/\|\hat y(\mathrm{C0})\|_2$ over the 32 grid times (also with weights and inputs cast to float64); the loss contrast $r_s=(\overline L_{\mathrm{native}}-\overline L_{\mathrm{resampled}})/\overline L_{\mathrm{resampled}}$ of the per-trajectory MSE; and wall time.
Uncertainty is a paired bootstrap over the 256 test trajectories with the seed fixed, with ratios recomputed in every draw.
Every interval is pointwise, for one estimand and one seed; it is not adjusted for the number of secondary conditions and seeds, and for a non-negative discrepancy an interval excluding zero shows that the change is nonzero, not what causes it.
Because the seeds share the trajectories (a crossed design), we report the seeds separately and give only their range, not a pooled interval.
A loss difference is called negligible only if the interval of $r_s$ lies inside $\pm5\%$, a convention (about 2.5\% in RMSE) rather than an application tolerance; non-significance is never read as equivalence.

% AUTO-GENERATED by scripts/make_paper_assets.py from results/raw -- do not edit by hand.
\begin{table}[t]
\caption{Model-coupling effect (P1-REAL-02, training seeds 1--4; same task, data and checkpoint rule; shared test set of 256 trajectories, ids 10000--10255).
Step and calib.\ MSE: selected checkpoint and its calibration MSE (label units squared).
S8 disc.: mean over trajectories of the relative $\ell_2$ change of the final per-step outputs at the 32 training-grid times when every held interval is split into 8 sub-steps, against C0 (float32).
Brackets: 95\% paired bootstrap CI over test trajectories with the seed fixed.}
\label{tab:seeds-coupling}
\centering\small
\begin{tabular}{rccc}
\toprule
seed & step & calib.\ MSE & S8 disc.\ [95\% CI] \\
\midrule
1 & 1300 & $0.0127$ & $0.0113$ $[0.0106,\,0.0121]$ \\
2 & 1600 & $0.0135$ & $0.0219$ $[0.0203,\,0.0236]$ \\
3 & 800 & $0.0168$ & $0.0043$ $[0.0041,\,0.0045]$ \\
4 & 1800 & $0.0125$ & $0.0147$ $[0.0140,\,0.0153]$ \\
\bottomrule
\end{tabular}
\end{table}

% AUTO-GENERATED by scripts/make_paper_assets.py from results/raw -- do not edit by hand.
\begin{table*}[t]
\caption{Readout decomposition under the same-path density change H8 (first half of $[0,T]$ split by 8; 144 tokens), training seeds 1--4, label units, means over the 256 test trajectories.
$D_s$ (primary, chosen after the development run): mean of $|p^{\mathrm{mean}}(\mathrm{H8})-p(\mathrm{C0})|-|p^{\mathrm{time}}(\mathrm{H8})-p(\mathrm{C0})|$.
Per trajectory, exactly, $p^{\mathrm{mean}}(\mathrm{H8})-p(\mathrm{C0})=W+M$ and $p^{\mathrm{time}}(\mathrm{H8})-p(\mathrm{C0})=Q$ with $W=\sum_j(a_j-w_j)b_j=\tfrac{7}{18}(A-B)$, $M=\sum_ja_j(z_j-b_j)$, $Q=\sum_jw_j(z_j-b_j)$; $b_j$ is the C0 output of token $j$'s base interval, $z_j$ the native H8 output, $a_j=1/144$, $w_j=dt_j/T$, and $A$, $B$ are the mean C0 outputs over the first and second half of $[0,T]$.
$W$ uses only C0 outputs: it is the same for any model whose H8 outputs repeat the C0 output within each base interval.
The signed terms are not an additive decomposition of the absolute-value estimand $D_s$.
Brackets: 95\% paired bootstrap CI (seed fixed).}
\label{tab:seeds-readout}
\centering\small
\begin{tabular}{rcccccc}
\toprule
seed & $D_s$ [95\% CI] & $p^{\mathrm{mean}}-p(\mathrm{C0})$ & $W$ & $M$ [95\% CI] & $Q$ [95\% CI] & $|Q|$ \\
\midrule
1 & $0.578$ $[0.512,\,0.649]$ & $+0.462$ & $+0.476$ & $-0.014$ $[-0.018,\,-0.010]$ & $-0.012$ $[-0.015,\,-0.009]$ & $0.019$ \\
2 & $0.584$ $[0.517,\,0.655]$ & $+0.473$ & $+0.488$ & $-0.016$ $[-0.019,\,-0.012]$ & $-0.015$ $[-0.017,\,-0.012]$ & $0.023$ \\
3 & $0.602$ $[0.536,\,0.673]$ & $+0.488$ & $+0.491$ & $-0.003$ $[-0.007,\,+0.001]$ & $-0.001$ $[-0.003,\,+0.001]$ & $0.013$ \\
4 & $0.602$ $[0.536,\,0.672]$ & $+0.497$ & $+0.486$ & $+0.011$ $[+0.008,\,+0.015]$ & $+0.014$ $[+0.012,\,+0.016]$ & $0.018$ \\
\bottomrule
\end{tabular}
\end{table*}

% AUTO-GENERATED by scripts/make_paper_assets.py from results/raw -- do not edit by hand.
\begin{table*}[t]
\caption{Task-risk effect: loss contrast $r_s=(\overline{L}_{\mathrm{native}}-\overline{L}_{\mathrm{resampled}})/\overline{L}_{\mathrm{resampled}}$ in percent (per-trajectory MSE at the 32 training-grid times; numerator and denominator recomputed in every bootstrap draw), training seeds 1--4.
Verdict: $=$ CI inside $\pm5\%$; $<$ / $>$ native lower / higher (CI excludes 0); $\ll$ / $\gg$ also beyond the margin; ? inconclusive. The $\pm5\%$ margin is a convention.
S2, S8 and H8 keep the held path; J1 (random observation times) changes the observed path. Timing: \cref{tab:timing}.}
\label{tab:seeds-loss}
\centering\small
\begin{tabular}{rcccccccc}
\toprule
seed & \multicolumn{2}{c}{$r_s$(S2) [95\% CI]} & \multicolumn{2}{c}{$r_s$(S8) [95\% CI]} & \multicolumn{2}{c}{$r_s$(H8) [95\% CI]} & \multicolumn{2}{c}{$r_s$(J1) [95\% CI]} \\
\midrule
1 & $-4.3$ $[-5.8,\,-2.7]$ & $<$ & $-6.2$ $[-8.8,\,-3.4]$ & $<$ & $-4.9$ $[-7.2,\,-2.3]$ & $<$ & $-21.9$ $[-32.0,\,-8.8]$ & $\ll$ \\
2 & $-2.1$ $[-5.3,\,+1.3]$ & ? & $+3.1$ $[-2.5,\,+9.4]$ & ? & $+6.7$ $[+1.8,\,+11.9]$ & $>$ & $-21.2$ $[-31.7,\,-6.5]$ & $\ll$ \\
3 & $-1.2$ $[-2.4,\,-0.2]$ & $=$ & $-1.9$ $[-3.9,\,-0.1]$ & $=$ & $-0.6$ $[-1.7,\,+0.6]$ & $=$ & $-21.8$ $[-31.9,\,-8.9]$ & $\ll$ \\
4 & $-5.4$ $[-8.0,\,-2.8]$ & $<$ & $-4.5$ $[-9.6,\,+0.7]$ & ? & $-2.1$ $[-5.7,\,+1.9]$ & ? & $-22.6$ $[-32.9,\,-9.0]$ & $\ll$ \\
\bottomrule
\end{tabular}
\end{table*}


\paragraph{Model coupling.}
The S8 discrepancy is $\RtwoSeightMin$--$\RtwoSeightMax$ in the four seeds, each interval excluding zero (\cref{tab:seeds-coupling}); it varies $\RtwoSeightFold$-fold between seeds (the earliest-selected checkpoint, step 800, has the smallest value, a relation not tested).
Casting weights and inputs to float64 changes the mean by at most $\RtwoFpEstMax$ and any per-trajectory value by at most $\RtwoFpTrajMax$, far below the smallest per-trajectory discrepancy $\RtwoDiscTrajMin$, so this is not float32 rounding; the pre-registered $10^{-4}$ threshold came from a scalar emulation (\cref{app:numerics}) and is not a floor of a trained 256-token network.
In the development checkpoint the change is absent from the encoder and first block ($\le\RealLayerOneFp$ in float32, $\RealLayerOneDp$ in float64) and enters in the second block ($\RealLayerTwoMin$--$\RealLayerTwoMax$), qualitatively consistent with \cref{prop:split,prop:cascade}; this is exploratory and rests on one checkpoint (\cref{app:pilot}).

\paragraph{Pooled weighting.}
$D_s$ is positive with an interval excluding zero in all four seeds ($\RtwoDmin$--$\RtwoDmax$, smallest lower bound $\RtwoDlomin$; \cref{tab:seeds-readout}), so the pre-specified replication criterion is met, but the cause is mostly the readout's weights: per trajectory the token-mean change splits exactly into $W=\tfrac{7}{18}(A-B)$, computed from the unchanged C0 outputs (\cref{sec:readout-analysis}), and a model term $M$.
The change averages $\DecDeltaMin$ to $\DecDeltaMax$, $W$ averages $\DecWMin$ to $\DecWMax$ and $M$ averages $\DecMMin$ to $\DecMMax$; $W$ has the sign of the change in at least $\DecSignMin\%$ of trajectories, and a token mean of outputs that merely repeat C0 would err by $\DecRepErrMin$--$\DecRepErrMax$ against the pooled label, as much as the native token mean ($\RtwoPoolMeanMin$--$\RtwoPoolMeanMax$).
The effect is therefore sample-count pooling of a non-constant output, which any per-token model shows, not a failure specific to selective SSMs.
$|M|$ averages $\DecMabsMin$--$\DecMabsMax$, and $\DecMendMin$--$\DecMendMax$ on the tokens that end a base interval, where it is the change of the model output itself.
Time weighting reduces the mean absolute change against C0 from $\DecDeltaAbsMin$--$\DecDeltaAbsMax$ to $\DecQabsMin$--$\DecQabsMax$ but does not remove it, and its pooled error ($\RtwoPoolTimeMin$--$\RtwoPoolTimeMax$) stays close to the resampled arm's ($\RtwoPoolResMin$--$\RtwoPoolResMax$).
$W$, $M$ and $Q$ come from one frozen forward pass per checkpoint that reproduces the stored pooled values exactly (\cref{app:prereg}); the signed terms do not decompose the absolute-value estimand $D_s$.

\paragraph{Task risk: the accuracy consequence does not replicate.}
The development run had suggested slightly better native accuracy on refined grids.
Across seeds 1--4, $r_s$ ranges from $\RtwoRSeightMin\%$ to $\RtwoRSeightMax\%$ on S8 and from $\RtwoRHeightMin\%$ to $\RtwoRHeightMax\%$ on H8 (\cref{tab:seeds-loss}): seed~2 is worse on H8 with an interval that excludes zero, seed~3 shows no practical difference on every held-path condition, and the rest are mixed, so we report no stable difference.
On J1 native MSE is lower in every seed ($r_s$ from $\RtwoRJoneMin\%$ to $\RtwoRJoneMax\%$), but the development checkpoint gives $\RtwoRJoneDev\%$ here and $\RtwoRJonePilotOld\%$ (inconclusive) on its own test set, so the J1 contrast depends on the test draw; J1 is not a refinement of the held path.

\paragraph{Cost.}
With the pre-specified protocol (inference mode, two warm-up and five timed calls, one batch of 256 sequences, 2 threads; \cref{app:real}), native\_S8\_end\_to\_end\,/\,resampled\_S8\_end\_to\_end is $\RtwoRatioMin$--$\RtwoRatioMax$ across seeds, forward-only $\RtwoFratioMin$--$\RtwoFratioMax$, and on C0 $\RtwoCzeroRatioMin$--$\RtwoCzeroRatioMax$ (\cref{tab:timing}); end-to-end includes resampling ($\RtwoResOverMin$--$\RtwoResOverMax$\,s on S8).
The two boundaries are timed in separate jobs and cross in $\RtwoInvN$ of $\RtwoInvTot$ cells (by up to $\RtwoInvMaxPct$\%), so differences of a few percent are run-to-run variation.

\paragraph{What this does not show.}
Four seeds of one two-block model on one synthetic task with a teacher from the model's own family replicate within that task only (\cref{sec:limits}); the primary estimand was chosen after the development run, and the calibration curves were unstable in every seed (up to $\RealDevMax$ in the development run), so the selected checkpoints depend on the evaluation schedule.

\subsection{Real sensor windows: UCI HAR, one checkpoint}
\label{sec:har}

\paragraph{Data and representation.}
We use the Human Activity Recognition Using Smartphones data~\citep{anguita2013har}: 30 subjects, six activities, 3-axis total acceleration, body acceleration and gyroscope at 50\,Hz, distributed as noise-filtered windows of 128 readings (2.56\,s) with 50\% overlap between consecutive windows of a subject; we use the nine inertial channels, not the 561 engineered features.
These distributor-preprocessed windows are not a raw stream, and the refinement below is artificial.
The official subject partition is kept (21 training subjects, 9 test subjects, $\HarNwin$ test windows); four training subjects, chosen before the run by a hash of their identifiers, are the calibration subjects ($\HarNcal$ windows), and the other 17 ($\HarNfit$ windows) fit the model and the per-channel normalization, which is never re-fitted.
A window's 128 values are represented as 128 held intervals of $1/50$\,s ($T{=}2.56$\,s), a chosen reconstruction: C0 keeps them; S8 splits every interval into 8 (1024 tokens); H8 splits the first 64 (576 tokens); held values, horizon and label are unchanged.
Because windows overlap, every quantity is averaged within a subject first and the 9 subject values are averaged with equal weight.

\paragraph{Model, training and arms.}
The model is the official \texttt{TIDESClassifier} of the pinned code with the backbone setting of \cref{sec:real-syn} ($\HarParams$ parameters), nine input channels and the official head, a token mean of the features followed by a linear map to six classes.
One model seed; Adam, batch 32, $\HarUpdatesRun$ updates; the state with the lowest calibration cross-entropy (every 100 updates) is frozen and hashed before any test window is read: update $\HarUpdate$, calibration cross-entropy $\HarCalCE$, accuracy $\HarCalAcc$.
The same checkpoint is evaluated with three arms: the official token-mean readout on the condition's grid (\emph{native mean}); the same features pooled with weights $dt_j/T$ at the same location and passed through the same affine head (\emph{native time}, an inference-time ablation, not a trained baseline); and the official readout on the training grid rebuilt from the observed values (\emph{resampled}), which on S8 and H8 is the identity control (input difference $\HarResIdent$).
Because the head $A(h)=Wh+b$ is affine and the weights sum to one, $A(\sum_jw_jh_j)=\sum_jw_jA(h_j)$ for both poolings (checked, \cref{app:har}); softmax outputs are never averaged.
The primary estimand, fixed before the run, is $d_s$, the mean over the windows of test subject $s$ of CE(native mean, H8) $-$ CE(native time, H8), and $D$, its equal-weight mean over the 9 subjects, with a subject-level paired percentile bootstrap (2000 draws); cross-entropy, accuracy, prediction flips, output distances and cost are secondary and reported separately.

% AUTO-GENERATED by scripts/make_paper_assets.py from results/raw -- do not edit by hand.
\begin{table*}[t]
\caption{P1-HAR-01: one official TIDES classifier (one model seed) on UCI HAR inertial windows, per official test subject (288--381 windows each; windows overlap by 50\%, so the subject is the analysis unit). C0: the 128 distributed values as 128 held intervals of $1/50$\,s; H8: the first 64 intervals split by 8 (576 tokens), same held values and labels. CE: mean cross-entropy (nats) of the window label under token-mean pooling (official head) and time-weighted pooling of the same features (inference ablation). $d_s$: primary difference CE(mean, H8) $-$ CE(time, H8) per subject; the last row is the equal-weight mean over the 9 subjects with a subject-level paired percentile bootstrap interval (2000 draws), which is coarse with 9 clusters. acc.: accuracy (time-weighted pooling: $0.903$ on H8). flip: fraction of windows whose predicted class under H8 (token mean) differs from C0.}
\label{tab:har-subjects}
\centering\small
\setlength{\tabcolsep}{5pt}
\begin{tabular}{rccccccc}
\toprule
subject & CE C0 & CE H8 mean & CE H8 time & $d_s$ & acc.\ C0 & acc.\ H8 mean & flip H8 mean \\
\midrule
2 & $0.193$ & $0.217$ & $0.194$ & $+0.023$ & $0.891$ & $0.884$ & $0.010$ \\
4 & $0.244$ & $0.220$ & $0.264$ & $-0.044$ & $0.896$ & $0.909$ & $0.050$ \\
9 & $0.455$ & $0.469$ & $0.432$ & $+0.038$ & $0.774$ & $0.771$ & $0.094$ \\
10 & $1.132$ & $1.103$ & $1.083$ & $+0.020$ & $0.714$ & $0.714$ & $0.054$ \\
12 & $0.106$ & $0.110$ & $0.100$ & $+0.011$ & $0.991$ & $0.994$ & $0.009$ \\
13 & $0.119$ & $0.144$ & $0.132$ & $+0.012$ & $0.954$ & $0.954$ & $0.012$ \\
18 & $0.072$ & $0.087$ & $0.075$ & $+0.013$ & $0.962$ & $0.962$ & $0.011$ \\
20 & $0.075$ & $0.140$ & $0.078$ & $+0.063$ & $0.952$ & $0.955$ & $0.008$ \\
24 & $0.037$ & $0.049$ & $0.038$ & $+0.011$ & $1.000$ & $1.000$ & $0.000$ \\
\midrule
mean & $0.270$ & $0.282$ & $0.266$ & $+0.016$ $[-0.004,\,+0.032]$ & $0.904$ & $0.905$ & $0.028$ \\
\bottomrule
\end{tabular}
\end{table*}

% AUTO-GENERATED by scripts/make_paper_assets.py from results/raw -- do not edit by hand.
\begin{table*}[t]
\caption{P1-HAR-01 effects of the same-held-path refinement, equal-weight means over the 9 test subjects (subject means first). $\Delta$CE: cross-entropy of the arm on the condition minus cross-entropy of the official token-mean readout on C0, with the subject-level bootstrap interval. flip: fraction of windows whose predicted class differs from C0. TV: total variation $\tfrac12\sum_c|p_c-p_c^{\mathrm{C0}}|$ between the class probabilities. c.-logit: $\|\tilde z-\tilde z^{\mathrm{C0}}\|_2/\|\tilde z^{\mathrm{C0}}\|_2$ with class-centered logits (0 windows excluded by the $10^{-6}$ rule). The resampled arm reproduces the C0 input exactly (identity control). On S8 the grid is uniform, so the two native arms coincide.}
\label{tab:har-effects}
\centering\small
\begin{tabular}{llcccc}
\toprule
cond. & arm & $\Delta$CE vs.\ C0 [95\% CI] & flip & TV & c.-logit \\
\midrule
S8 & native, token mean & $-0.0097$ $[-0.034,\,+0.012]$ & $0.008$ & $0.008$ & $0.029$ \\
S8 & native, time-weighted & $-0.0097$ $[-0.034,\,+0.012]$ & $0.008$ & $0.008$ & $0.029$ \\
S8 & resampled (identity control) & $+0.0000$ $[+0.000,\,+0.000]$ & $0.000$ & $0.000$ & $0.000$ \\
\midrule
H8 & native, token mean & $+0.0119$ $[-0.005,\,+0.028]$ & $0.028$ & $0.037$ & $0.210$ \\
H8 & native, time-weighted & $-0.0044$ $[-0.017,\,+0.007]$ & $0.003$ & $0.005$ & $0.016$ \\
H8 & resampled (identity control) & $+0.0000$ $[+0.000,\,+0.000]$ & $0.000$ & $0.000$ & $0.000$ \\
\bottomrule
\end{tabular}
\end{table*}


\paragraph{Results (\cref{tab:har-subjects,tab:har-effects}).}
\emph{Pooled splitting change and its common-time part.}
Splitting every held interval (S8) changes the class-centered pooled logits by $\HarCLSeightMean$ (relative $\ell_2$; subject range in \cref{tab:har-subjects}), the class probabilities by a total variation of $\HarTVSeightMean$, and the predicted class of $\HarFlipSeightMean$ of the windows, while the resampled arm reproduces C0 exactly; the cross-entropy change is $\HarDCESeightMean$ nats with a subject-level interval of $[\HarDCESeightMeanlo,\HarDCESeightMeanhi]$, so no risk change is established at nine subjects.
Unlike the synthetic discrepancy of \cref{sec:real-syn}, which compares outputs at the same 32 times, this pooled change compares a mean over 128 time points with a mean over 1024, so it is not by itself a common-time (model-coupling) effect: even an exact single layer has a split-invariant end value but a different token mean.
A fixed-checkpoint diagnostic designed after these results (\cref{tab:har-ct}, \cref{app:har}) separates the two with the official affine head $A$: $z_0=A(\overline{H_0})$, $z_{\mathrm{end}}=A(\mathrm{mean}_kH_8[\mathrm{end}(k)])$ from the S8 features at the 128 common time points, $z_{\mathrm{all}}=A(\overline{H_8})$, and exactly $z_{\mathrm{all}}-z_0=(z_{\mathrm{end}}-z_0)+(z_{\mathrm{all}}-z_{\mathrm{end}})$.
The common-time change has relative norm $\HarCtRelEnd$ and the extra-sampling change $\HarCtRelExtra$ (per-window ratios $\|\Pi d\|_2/\|\Pi z_0\|_2$ averaged subject-first, as in \cref{tab:har-effects}; as a ratio of subject-mean norms the common-time part is $\HarCtRatioOfMeans$), with inner product $\HarCtInner$ logit$^2$, so the two norms are not additive shares; the final features at the common end points move by $\HarCtFeatEnd$ relative, and $z_{\mathrm{end}}$ flips no prediction (cross-entropy $\HarCtCEend$ against $\HarCtCEzero$ for $z_0$).
The model's outputs at the common times therefore do change under splitting, but by an order of magnitude less than the pooled readout; the extra-sampling term involves the model path through the added tokens and is not a pure quadrature error, and neither term is attributed to a layer.
\emph{Pooled weighting.}
Under H8 the token-mean readout moves the centered logits by $\HarCLHeightMean$, changes the predicted class of $\HarFlipHeightMean$ of the windows and the cross-entropy by $\HarDCEHeightMean$ $[\HarDCEHeightMeanlo,\HarDCEHeightMeanhi]$; time weighting of the same features reduces these to $\HarCLHeightTime$, $\HarFlipHeightTime$ and $\HarDCEHeightTime$ $[\HarDCEHeightTimelo,\HarDCEHeightTimehi]$.
The logit decomposition of \cref{sec:real-syn} applies unchanged (norms in logit units, per window, then subject means): the token-mean change has norm $\HarDecDelta$, the density-weighting term from the C0 token logits alone $\HarDecW$, the model term $M$ $\HarDecM$, its part on the 128 interval-end tokens $\HarDecMend$, and the time-weighted change $Q$ $\HarDecQ$: as on the synthetic task, the large token-mean effect is mostly the weighted-average identity, and the model outputs at the base times change two orders of magnitude less.
\emph{Task risk.}
The primary estimand is $D=\HarD$ nats with a subject-level interval of $[\HarDlo,\HarDhi]$; $d_s$ is positive for $\HarDpos$ of the 9 subjects (range $\HarDmin$ to $\HarDmax$).
The interval includes zero: the pre-registered difference is not established on this dataset, and non-significance is not evidence of equivalence; the sign pattern matches the synthetic finding, but nine clusters give a coarse interval.
Accuracy is $\HarAccCzero$ on C0, $\HarAccHeightMean$ under H8 with the token mean and $\HarAccHeightTime$ with time weighting (no interval is reported).
\emph{Control and cost.}
A control (time-weighted per-channel means and second moments with a multinomial linear classifier; not a task ceiling or a capacity-matched baseline) has identical features on C0, S8 and H8 (maximum difference $\HarCtlFeat$; cross-entropy $\HarCtlCE$, accuracy $\HarCtlAcc$), as any readout that integrates the held path over time must.
Native evaluation costs native\_S8\_end\_to\_end\,/\,resampled\_S8\_end\_to\_end $=\HarRatioSeight$ (forward-only $\HarFratioSeight$) and $\HarRatioHeight$ on H8 (\cref{tab:har-timing}).

\paragraph{Initialization sensitivity (two further model seeds).}
The frozen protocol was retrained with model seeds 101 and 102 and each checkpoint evaluated once on the \emph{same} nine subjects (\cref{app:har}): a sensitivity check on one architecture and one dataset, not an independent replication; the three runs are separate rows, never 27 observations, and the pre-registered primary of the original run stands.
The separation repeats: the common-time part of the pooled S8 change is $\RviiiCommonA$ and $\RviiiCommonB$ against pooled changes of $\RviiiCLSeightA$ and $\RviiiCLSeightB$, $z_{\mathrm{end}}$ flips no prediction, and $W$ dominates the H8 token-mean change in each run (\cref{tab:har-r8-decomp}).
The risk numbers move more: $D=\RviiiDA$ $[\RviiiDAlo,\RviiiDAhi]$ ($\RviiiDApos$ of 9 subjects positive) and $\RviiiDB$ $[\RviiiDBlo,\RviiiDBhi]$ ($\RviiiDBpos$ of 9) against $\HarD$, with H8 token-mean cross-entropy changes against C0 of $\RviiiDCEHeightMeanA$ and $\RviiiDCEHeightMeanB$ (time-weighted $\RviiiDCEHeightTimeA$, $\RviiiDCEHeightTimeB$): the direction is consistent across the three initializations, the size varies about three-fold, and the two new intervals exclude zero but share subjects and data with the original run and do not change its verdict.
\paragraph{What this does not show.}
One dataset and one architecture give results conditional on these checkpoints, not on the population of training runs, other sensors, naturally irregular acquisitions or TIDES' benchmarks; nine overlapping-window subjects bound what the interval can resolve, the common-time diagnostic is post hoc on the original checkpoint, and the two further seeds reuse its subjects and data.

%=====================================================================================
\section{Controlled Toy Results}
\label{sec:toy}

These fixed-parameter results are numerical checks and controlled measurements, not proofs, and say nothing directly about trained models.
The pre-registered FIRST\_RUN used \NunitsPrimary{} independent units (one parameter and input draw each), $T{=}8$, $K{=}16$, $m\in\{1,2,4,8\}$, 2 CPU threads and \FirstRunSeconds\,s; details and deviations are in \cref{app:prereg}.

\paragraph{Splitting.}
With $\Delta{=}dt\,g$ and exact ZOH the output is split-invariant to $\EoneZohMax$ and agrees with an independent RK4 reference to $\EoneRKmax$ (an engineering check of \cref{prop:split}).
Euler-$B$ with the correct step converges at first order (median slope $\EoneEulerSlope$) but has a median base-step artifact of $\EoneEulerMone$ in this toy; the magnitude depends on the toy's $\Delta|\lambda|$ scale.
Bilinear converges at second order ($\EoneBilinSlope$); with $\Delta{=}\tau g$ the artifact grows from $\EoneNodtMtwo$ to $\EoneNodtMeight$ (medians, $m{=}2\to8$), by construction (\cref{tab:split}, \cref{fig:convergence}a).

% AUTO-GENERATED by scripts/make_paper_assets.py from results/raw -- do not edit by hand.
\begin{figure}[t]
\centering
\definecolor{seriesA}{HTML}{2A78D6}\definecolor{seriesB}{HTML}{EB6834}\definecolor{seriesC}{HTML}{1BAF7A}
\begin{tikzpicture}
\begin{groupplot}[group style={group size=2 by 1, horizontal sep=0.85cm},
  width=0.47\columnwidth, height=0.47\columnwidth, xmode=log, ymode=log, log basis x=2,
  xtick={1,2,4,8}, xticklabels={1,2,4,8}, xlabel={sub-steps $m$}, tick label style={font=\footnotesize},
  label style={font=\footnotesize}, title style={font=\footnotesize}, grid=major, grid style={gray!20, line width=0.3pt},
  axis line style={gray!60}, legend style={font=\scriptsize, draw=none, fill=none}, legend cell align=left]
\nextgroupplot[title={(a) single layer}, ylabel={relative error}, ymin=5e-5, ymax=5,
  legend style={at={(0.5,-0.33)}, anchor=north}]
\addplot[seriesA, line width=1pt, mark=*, mark size=1.6pt] coordinates {(1,2.103473e-01) (2,1.022642e-01) (4,5.040279e-02) (8,2.495612e-02)};
\addlegendentry{Euler-$B$, $\Delta{=}dt\,g$}
\addplot[seriesB, line width=1pt, dashed, mark=square*, mark size=1.6pt] coordinates {(1,9.454201e-03) (2,2.284158e-03) (4,5.578572e-04) (8,1.386879e-04)};
\addlegendentry{bilinear, $\Delta{=}dt\,g$}
\addplot[seriesC, line width=1pt, dotted, mark=triangle*, mark size=2pt] coordinates {(2,4.762033e-01) (4,9.339680e-01) (8,1.275009e+00)};
\addlegendentry{exact ZOH, $\Delta{=}\tau g$}
\nextgroupplot[title={(b) two-layer cascade}, ymin=5e-3, ymax=0.5,
  legend style={at={(0.5,-0.33)}, anchor=north}]
\addplot[seriesA, line width=1pt, mark=*, mark size=1.6pt] coordinates {(1,1.884736e-01) (2,1.001190e-01) (4,5.162333e-02) (8,2.621494e-02)};
\addlegendentry{$a{=}1,\,c{=}0.5$}
\addplot[seriesB, line width=1pt, dashed, mark=square*, mark size=1.6pt] coordinates {(1,1.966501e-01) (2,1.063295e-01) (4,5.535558e-02) (8,2.825132e-02)};
\addlegendentry{$a{=}c{=}1$}
\end{groupplot}
\end{tikzpicture}
\caption{Error against the exact held-path solution as the same held interval is split into $m$ sub-steps: (a) single layer (FR-E1-SPLIT; medians over 32 toy units); (b) two-layer cascade (P1-COMP-01; one deterministic case per line).
Exact ZOH with $\Delta{=}dt\,g$ (a) and the exact coupled solution (b) stay at $10^{-15}$ and are not drawn. The $\Delta{=}\tau g$ line starts at $m{=}2$ because it is exact at $m{=}1$ by construction.
Values are listed in \cref{tab:split,tab:coupling}.}
\label{fig:convergence}
\end{figure}


\paragraph{New observations.}
With exact ZOH the artifact against RK4 on the new held path stays below $\EtwoZohArt$, so the whole change is a response to the changed input reconstruction.
The exact decomposition \eqref{eq:decomp} of the other variants (exploratory, \cref{app:toy}) has cross terms between $\EtwoCrossMin$ and $\EtwoCrossMax$ of the per-unit total, so no per-unit share is defined.

\paragraph{Readouts and interaction of fixes.}
On a split-invariant backbone, a density change moves the sample mean by $\EthreeMeanChange$ and the exact integral by $\EthreeExactChange$ (scale-normalized, $m{=}8$).
The pre-registered clause H7c was mis-specified before the run and failed (slope $\HsevenOrigSlope$); its restatement passed only weakly (median slope $\EthreeRMEslope$, $\EthreeRMEneg$ of 32 units negative).
One fix at a time (\cref{tab:interaction}): exact ZOH alone gives $\FacZoh$, the physical step alone $\FacDt$, a time readout alone $\FacTime$, against $\FacNone$ without fixes; in this tested ablation the three corrections combined reduced the artifact to $\FacAll$.
This is not a necessity result: it concerns a time-average readout of one toy, a last-state readout needs no time integral, and special inputs or coefficients (for example a zero input) remove it without all three corrections.
Under the pre-registered stop condition (an operational rule, not a statistical, equivalence or novelty test) the toy gives no support for a new architecture.

\paragraph{Two-layer cascade.}
For \eqref{eq:cascade} (P1-COMP-01; held input of FIRST\_RUN unit~0; \CompSeconds\,s), the exact coupled solution is split-invariant to $\CompInv$, while the layerwise computation has error $\CompErrOne$ at the base step and $\CompErrEight$ at $m{=}8$ (slope $\CompSlope$; $\CompSlopeEq$ for $a=c$; \cref{fig:convergence}b, \cref{tab:coupling}), consistent with, not a proof of, first order for this input.


%=====================================================================================
\section{Limitations}
\label{sec:limits}

\emph{Synthetic evidence:} four training seeds and a development seed of one small two-block model (one selective SSM) on one synthetic task whose teacher is from the model's own family; it supports no statement about other tasks, the population of training runs, TIDES' benchmarks, Mamba, S4 or S5, and the accuracy contrast was not stable. \emph{Real-sensor evidence:} one dataset, one architecture, three model seeds on the same nine subjects (one pre-registered run and a two-seed sensitivity check), an artificial refinement of distributor-preprocessed windows (\cref{sec:har}); the task-risk effect keeps its direction but its size varies across seeds, and no other model or dataset was evaluated.
\emph{Entry points not measured:} token-indexed convolution, normalization in training mode, and deeper or wider stacks.
\emph{Analysis:} \cref{prop:split,prop:cascade} are standard; \cref{prop:cascade} is an upper bound for a linear time-invariant scalar cascade (\cref{rem:scope}), with no lower bound, no statement about deep selective networks and no global hold-error or Euler-$B$ bound.
\emph{Post hoc choices:} the primary estimand of the repetition was chosen after the development run; the layer localization, the H8 signed decomposition, the HAR common-time diagnostic, the toy decomposition and the stiffness columns were analysed after the results were seen.

%=====================================================================================
\section{Conclusion}
\label{sec:conclusion}

Grid dependence in selective SSMs can enter through a step that counts samples, an inexact input-matrix discretization, sample-count readouts and, in stacks, the resampling of internal signals; the last is not caused by any change of the observed input.
In the single-layer toy ablation only the combined corrections reduced the readout artifact to rounding level, and the inter-layer resampling error remained ($O(\text{mesh})$ upper bound for a scalar linear cascade).
On the official TIDES code (physical step), splitting the held path changes the outputs at common times in all four training seeds of one synthetic task; on real sensor windows (UCI HAR) the common-time part of the pooled splitting change is 8- to 16-fold smaller than the readout term over the added tokens (not additive shares) in all three model seeds; the larger readout effect of a same-path density change is, in both, mostly the weighted-average identity of token-mean pooling, which time weighting reduces to a small residual; and the loss consequence does not replicate across synthetic seeds, while on the sensor windows the pre-registered subject-level interval includes zero and two further seeds give larger positive values on the same subjects.
Whether these effects matter for other selective SSMs, for naturally irregular acquisitions, and across training runs is open.\label{body-end}

%=====================================================================================
\section*{Impact Statement}

This paper presents work whose goal is to advance the field of Machine Learning, specifically the understanding of how sequence models respond to changes in sampling.
The experiments use synthetic data and one public, anonymized activity-recognition dataset (UCI HAR, distributed for research with a citation requirement), all on small CPU computations; the trained-model results come from four training seeds and a development seed on one synthetic task and from one checkpoint on the public dataset. No clinical or otherwise sensitive data were used.
The main risk we see is over-reading such results: a model that behaves consistently in a controlled setting may still behave differently on irregularly sampled real data, for example clinical or sensor streams, so any use in such settings needs evaluation on the actual data and sampling process.
We are not aware of other societal consequences that must be specifically highlighted here.

\bibliography{references}
\bibliographystyle{icml2026}

%=====================================================================================
\newpage
\appendix
\onecolumn

\FloatBarrier
\section{Pre-registration, deviations and compute}
\label{app:prereg}

\paragraph{FIRST\_RUN.}
Hypotheses H1--H8, metrics, 32 units, 2 CPU threads, a 120\,s budget and stop conditions were committed before any experiment code ran.
Amendments recorded before any experiment ran: A1 (adaptive RK4 reference step), A2 (a restated H7c, keeping the mis-specified original), A3 (secondary configurations fixed to seeds 0--7, descriptive).
Deviations: a timing-only smoke test on seed 999 (metrics not inspected); a repeated execution after a manifest bookkeeping bug, with byte-identical data files; a partial execution killed by a shell pipe, whose outputs were deleted.
P1-COMP-01 was pre-registered separately before its code ran.
\Cref{tab:hypotheses} lists every pre-registered verdict, including the failure.
\emph{Literature:} the key equations and experiment sections of TIDES, Mamba-3 and S5 were read in the original text; other entries were checked by keyword search.

\paragraph{P1-REAL-01.}
The design, the pinned implementation and the predictions were committed before any installation.
Amendments R-A1--R-A9 were committed after an untrained adapter check and a timing smoke on four training trajectories, and before training: calibration-based checkpoint selection, same-held-path refinements only (new-observation conditions removed, one random-time condition kept as secondary), a resampler restricted to observed values, the pooled readout and pooling-only arm, the explicit equivalence unit, the single-seed pilot label, time caps, and seeding of all random generators (the TIDES initialization also draws from NumPy's global generator).
The pre-registered predictions were: native splitting discrepancy at S8 above $10^{-4}$ (supported), resampled splitting discrepancy exactly zero (supported), native runtime growing with tokens (observed), and a larger sample-mean than time-weighted pooled change on H8 (supported).

\paragraph{P1-REAL-02.}
The configuration (seeds, identical protocol, new test set, primary and secondary estimands, bootstrap and margin rules, timing boundaries, failure handling, conditional extra diagnostics) was committed before any repetition run.
One incident occurred: the evaluation of the development row crashed before writing any output because its output directory was not created; the attempt is recorded as INVALID, the directory creation was added, and only that row was rerun (the four training seeds were unaffected).
The pre-specified conditional layer diagnostic was not triggered.

\paragraph{H8 signed decomposition (derived, after the results).}
P1-REAL-02 stored per-trajectory pooled values but not per-token outputs.
A configuration fixed before execution specified one frozen native forward pass per existing checkpoint on C0 and H8 (no training, no timing, same test set), a check that it reproduces the stored pooled values (maximum absolute difference $\DecMismatch$), and the signed identities of \cref{tab:seeds-readout} (largest residual $\DecIdentity$).

\paragraph{Compute.}
All on CPU (2 threads for experiments), no GPU, no paid API; the synthetic stages use no external data, and the only external data are the public UCI HAR windows of \cref{app:har}; wall and process CPU seconds are distinguished where recorded.
FIRST\_RUN: two executions of about 21\,s, one partial execution and a smoke test of about 3\,s together.
P1-COMP-01: \CompSeconds\,s.
P1-REAL-01 (development): installation of a CPU tensor library (35\,s) and of a \LaTeX\ toolchain (66\,s); adapter check and timing smoke (${\approx}4$\,s; $\AdapterStep$\,s per batch-32 step); convention check (${\approx}2$\,s); training including label generation ($\RealTrainSec$\,s); evaluation ($\RealEvalSec$\,s); layer diagnostic (${\approx}11$\,s).
P1-REAL-02: data generation $\RtwoDataWall$\,s; seed-0 equivalence retrain ${\approx}57$\,s; four training runs $\RtwoTrainWall$\,s wall in total ($\RtwoTrainCPU$\,s process CPU in the training loops); test-set generation $\RtwoTestWall$\,s; five evaluations $\RtwoEvalWall$\,s wall in total, plus the invalid development attempt (${\approx}30$\,s); H8 decomposition forward passes $\DecWall$\,s.

% AUTO-GENERATED by scripts/make_paper_assets.py from results/raw -- do not edit by hand.
\begin{table}[h]
\caption{Pre-registered check verdicts, including the failed H7c\_original (mis-specified before the run; amendment A2) and the weakly supported H7c\_restated. These are bookkeeping, not scientific results.}
\label{tab:hypotheses}
\centering\small
\begin{tabular}{llc}
\toprule
run & check & verdict \\
\midrule
FIRST\_RUN & H1 & PASS \\
FIRST\_RUN & H2 & PASS \\
FIRST\_RUN & H3 & PASS \\
FIRST\_RUN & H4 & PASS \\
FIRST\_RUN & H5a & PASS \\
FIRST\_RUN & H5b & PASS \\
FIRST\_RUN & H6 & PASS \\
FIRST\_RUN & H7a & PASS \\
FIRST\_RUN & H7b & PASS \\
FIRST\_RUN & H7c\_original & FAIL \\
FIRST\_RUN & H7c\_restated & PASS (weak; see text) \\
FIRST\_RUN & H8 & PASS \\
P1-COMP-01 & COMP-H1 (distinct) & PASS \\
P1-COMP-01 & COMP-H2 (distinct) & PASS \\
P1-COMP-01 & COMP-H3 (distinct) & PASS \\
P1-COMP-01 & COMP-H4 (distinct) & PASS \\
P1-COMP-01 & COMP-H5 (distinct) & PASS \\
P1-COMP-01 & COMP-H1 (equal) & PASS \\
P1-COMP-01 & COMP-H2 (equal) & PASS \\
P1-COMP-01 & COMP-H3 (equal) & PASS \\
P1-COMP-01 & COMP-H4 (equal) & PASS \\
P1-COMP-01 & COMP-H5 (equal) & PASS \\
\bottomrule
\end{tabular}
\end{table}


\FloatBarrier
\section{Implementation and protocol details}
\label{app:real}

\emph{Pinned code facts (read, then executed through the adapter):} the step is \texttt{step\_scale * exp(log\_step)} with a per-mode timescale; the ZOH input matrix is \texttt{(1/Lambda)*(Lambda\_bar-1)*B} without \texttt{expm1}; the scan is $x_k=\bar\Lambda_kx_{k-1}+\bar B_ku_k$ with the gap before step $k$; each block applies \texttt{BatchNorm1d(affine=False)} over batch and time (running statistics at evaluation); an optional causal convolution defaults to off; the classifier head pools with \texttt{h.mean(dim=1)}.
\emph{Adapter check (untrained, four training trajectories):} forward outputs finite; gradients finite for all 54 parameters that receive one (the six static $\mathrm{Re}\,\Lambda$, $B$, $C$ tensors are unused in input-dependent mode; parameters upstream of the zero-initialized head projections receive zero gradient at step~0 only); an Adam step changed the parameters and lowered the loss on the same batch; evaluation was deterministic, and a per-step step tensor of ones matched the scalar step exactly.
\emph{Labels:} teacher output on the continuous path from an independent RK4 integration (128 sub-steps per training interval); the pooled label is its time average by Simpson's rule.
\emph{Readout of J1:} the last native output at or before each training-grid time.
\emph{Eval-mode checks (P1-REAL-02, every seed):} every module in evaluation mode during calibration and test; dropout probability 0; batch normalization with running statistics; equal sequence lengths within every batch (no padding or masking is used); all step sizes positive (the smallest native step is $1/8$ of the training step on S8 and H8).
\emph{Timing (P1-REAL-02):} \cref{tab:timing}.

% AUTO-GENERATED by scripts/make_paper_assets.py from results/raw -- do not edit by hand.
\begin{table}[h]
\caption{P1-REAL-02 timing (one batch of 256 sequences; median of 5 timed calls after 2 warm-ups; \texttt{inference\_mode}; 2 threads). End-to-end includes tensor construction, LOCF resampling for the resampled arm, forward, de-normalization and readout mapping; forward-only is the model call on prebuilt tensors. CPU seconds are process CPU time summed over threads. Timed repeats are not independent samples. $^\ddagger$Development checkpoint (seed 0).}
\label{tab:timing}
\centering\small
\begin{tabular}{rcrccccc}
\toprule
seed & cond. & tokens & nat.\ e2e & nat.\ fwd & res.\ e2e & res.\ fwd & nat.\ e2e CPU \\
\midrule
0$^\ddagger$ & C0 & 32 & $0.049$ & $0.044$ & $0.047$ & $0.042$ & $0.10$ \\
0$^\ddagger$ & S8 & 256 & $0.578$ & $0.574$ & $0.068$ & $0.063$ & $1.12$ \\
0$^\ddagger$ & H8 & 144 & $0.297$ & $0.281$ & $0.069$ & $0.056$ & $0.58$ \\
1 & C0 & 32 & $0.042$ & $0.039$ & $0.042$ & $0.041$ & $0.09$ \\
1 & S8 & 256 & $0.565$ & $0.545$ & $0.049$ & $0.041$ & $1.10$ \\
1 & H8 & 144 & $0.308$ & $0.275$ & $0.047$ & $0.042$ & $0.60$ \\
2 & C0 & 32 & $0.041$ & $0.038$ & $0.042$ & $0.038$ & $0.08$ \\
2 & S8 & 256 & $0.560$ & $0.550$ & $0.058$ & $0.039$ & $1.09$ \\
2 & H8 & 144 & $0.300$ & $0.281$ & $0.053$ & $0.043$ & $0.59$ \\
3 & C0 & 32 & $0.051$ & $0.044$ & $0.049$ & $0.043$ & $0.10$ \\
3 & S8 & 256 & $0.544$ & $0.566$ & $0.054$ & $0.047$ & $1.06$ \\
3 & H8 & 144 & $0.286$ & $0.275$ & $0.061$ & $0.047$ & $0.56$ \\
4 & C0 & 32 & $0.045$ & $0.041$ & $0.045$ & $0.042$ & $0.09$ \\
4 & S8 & 256 & $0.574$ & $0.531$ & $0.054$ & $0.046$ & $1.10$ \\
4 & H8 & 144 & $0.281$ & $0.276$ & $0.055$ & $0.047$ & $0.56$ \\
\bottomrule
\end{tabular}
\end{table}


\FloatBarrier
\section{Development pilot (seed 0) on its own test set}
\label{app:pilot}

\Cref{tab:dev-newtest} gives the development checkpoint on the new shared test set. The remaining results come from the development run (P1-REAL-01), whose questions and estimands were not fixed in advance beyond its own predictions, on its own test set (identifiers 640--895).
Its equivalence analysis used a margin of 5\% of the point-estimated resampled MSE (kept here as exploratory); the repetition in \cref{sec:real} instead bootstraps the ratio $r_s$.
Pooled quantities in \cref{tab:pool} are mean absolute errors (or changes) of pooled predictions against the pooled label, in label units; they are not hidden-state discrepancies or MSEs.
Its runtimes are forward-only model calls on prebuilt tensors (no warm-up, outside inference mode) and exclude resampling.
\emph{Layer diagnostic (exploratory, not pre-registered).} After seeing the development results we measured the change at intermediate features of the development checkpoint (\cref{tab:layers}): the encoder output is unchanged, the first block changes by at most $\RealLayerOneFp$ in float32 and $\RealLayerOneDp$ with the weights cast to float64, and the second block by $\RealLayerTwoMin$ to $\RealLayerTwoMax$ in both precisions. The first block receives the held input and uses an exact held-coefficient update, whereas the second block receives the first block's output, which splitting samples at more points within each interval; neither proposition is a statement about this trained network. The repetition pre-specified a per-seed layer diagnostic only if some seed's S8 discrepancy fell below $10^{-3}$; this did not happen, so the localization rests on one checkpoint.

% AUTO-GENERATED by scripts/make_paper_assets.py from results/raw -- do not edit by hand.
\begin{table}[h]
\caption{Development checkpoint (seed 0, P1-REAL-01) evaluated with the P1-REAL-02 protocol on the new shared test set; not part of the replication criterion. Definitions as in \cref{tab:seeds-coupling,tab:seeds-readout,tab:seeds-loss}.}
\label{tab:dev-newtest}
\centering\small
\begin{tabular}{ll}
\toprule
quantity & development checkpoint, new test set \\
\midrule
selected step / calibration MSE & 1500 / $0.0119$ \\
S8 disc.\ [95\% CI] & $0.0114$ $[0.0108,\,0.0120]$ \\
$D_s$ [95\% CI] & $0.575$ $[0.511,\,0.643]$ \\
$p^{\mathrm{mean}}-p(\mathrm{C0})$, $W$, $M$, $Q$ (means) & $+0.452$, $+0.460$, $-0.009$, $-0.009$ \\
$r_s$(S2) [95\% CI], verdict & $-2.9$ $[-4.6,\,-1.5]$ $=$ \\
$r_s$(S8) [95\% CI], verdict & $-3.9$ $[-6.7,\,-1.8]$ $<$ \\
$r_s$(H8) [95\% CI], verdict & $-6.5$ $[-9.4,\,-3.8]$ $<$ \\
$r_s$(J1) [95\% CI], verdict & $-20.9$ $[-30.8,\,-8.0]$ $\ll$ \\
\bottomrule
\end{tabular}
\end{table}

% AUTO-GENERATED by scripts/make_paper_assets.py from results/raw -- do not edit by hand.
\begin{table}[t]
\caption{DEVELOPMENT run (P1-REAL-01): trained official TIDES, seed 0, on its own test set (256 trajectories, ids 640--895; one checkpoint selected on the calibration split; paired).
``disc.'': mean relative $\ell_2$ change of the per-step outputs at the 32 training-grid times against the same arm on C0.
MSE: per-trajectory mean squared error to the teacher at those times (original units; label s.d.\ 2.35).
$\Delta$MSE: mean paired difference native $-$ resampled in units of $10^{-4}$ (label units squared) with a 95\% bootstrap CI; the margin is $\pm5\%$ of the point-estimated resampled mean MSE of the same condition (exploratory rule of the development run).
Verdict (native vs.\ resampled MSE): lower/higher if the CI excludes 0 on that side. Runtimes of this run are not tabulated; they were forward-only and are superseded by \cref{tab:timing}.
C0 is the training grid; S$m$ and H8 (first half split by 8) keep the held path; $^\dagger$J1 (random observation times) is a secondary non-refinement condition whose held path differs by construction.}
\label{tab:real}
\centering\small
\begin{tabular}{lcccccc}
\toprule
cond. & disc.\ native & disc.\ resampled & MSE native & MSE resampled & $\Delta$MSE [95\% CI] ($10^{-4}$) & margin rule \\
\midrule
C0 & $0$ & $0$ & $0.0167$ & $0.0167$ & $0$ $[0,\,0]$ & identical inputs (check) \\
S2 & $6.1\times10^{-3}$ & $0$ & $0.0162$ & $0.0167$ & $-5.1$ $[-7.3,\,-2.7]$ & within $\pm5\%$ \\
S4 & $9.2\times10^{-3}$ & $0$ & $0.0161$ & $0.0167$ & $-6.5$ $[-9.9,\,-2.9]$ & lower \\
S8 & $0.011$ & $0$ & $0.0160$ & $0.0167$ & $-6.9$ $[-10.8,\,-2.7]$ & lower \\
H8 & $7.8\times10^{-3}$ & $0$ & $0.0156$ & $0.0167$ & $-10.8$ $[-14.0,\,-7.6]$ & lower \\
J1$^\dagger$ & $0.35$ & $0.39$ & $1.14$ & $1.22$ & $-785$ $[-2120,\,+623]$ & inconclusive \\
\bottomrule
\end{tabular}
\end{table}

% AUTO-GENERATED by scripts/make_paper_assets.py from results/raw -- do not edit by hand.
\begin{table}[t]
\caption{Pooled readout of the development checkpoint on its own test set (P1-REAL-01), label units, means over 256 test trajectories. Left: $n^{-1}\sum_i|p_i(c)-\bar y_i|$; right: $n^{-1}\sum_i|p_i(c)-p_i(\mathrm{C0})|$, with $p=p^{\mathrm{mean}}$ (``native mean'', token mean as in the TIDES classifier head), $p=p^{\mathrm{time}}$ (``native time'', right-endpoint time-weighted sum of the same native outputs; a pooling-only ablation, not an exact integral) or $p=p^{\mathrm{res}}$ (``resampled'', token mean of the 32 resampled outputs); definitions in \cref{sec:real}. Pooled label $\bar y_i$: time average of the teacher output over $[0,T]$ by Simpson's rule on its RK4 nodes (mean absolute value 1.59).}
\label{tab:pool}
\centering\small
\begin{tabular}{lcccccc}
\toprule
 & \multicolumn{3}{c}{error to pooled label} & \multicolumn{3}{c}{change vs.\ C0} \\
\cmidrule(lr){2-4}\cmidrule(lr){5-7}
cond. & native mean & native time & resampled & native mean & native time & resampled \\
\midrule
C0 & $0.063$ & $0.063$ & $0.063$ & $0$ & $0$ & $0$ \\
S8 & $0.070$ & $0.070$ & $0.063$ & $0.032$ & $0.032$ & $0$ \\
H8 & $0.56$ & $0.068$ & $0.063$ & $0.55$ & $0.016$ & $0$ \\
J1 & $0.37$ & $0.28$ & $0.29$ & $0.36$ & $0.26$ & $0.30$ \\
\bottomrule
\end{tabular}
\end{table}

% AUTO-GENERATED by scripts/make_paper_assets.py from results/raw -- do not edit by hand.
\begin{table}[t]
\caption{Where the splitting change appears in the development checkpoint (EXPLORATORY, added after its results were seen). Mean relative $\ell_2$ change against C0 at the training-grid times of the
features after each block and of the (normalized) output, with the frozen weights evaluated in float32 and cast to float64. The encoder output is unchanged (0) in all cases.}
\label{tab:layers}
\centering\small
\begin{tabular}{lcccccc}
\toprule
 & \multicolumn{2}{c}{block 1} & \multicolumn{2}{c}{block 2} & \multicolumn{2}{c}{output} \\
\cmidrule(lr){2-3}\cmidrule(lr){4-5}\cmidrule(lr){6-7}
cond. & fp32 & fp64 & fp32 & fp64 & fp32 & fp64 \\
\midrule
S2 & $5.6\times10^{-6}$ & $5.1\times10^{-15}$ & $4.1\times10^{-3}$ & $4.1\times10^{-3}$ & $6.8\times10^{-3}$ & $6.8\times10^{-3}$ \\
S4 & $6.0\times10^{-6}$ & $1.1\times10^{-14}$ & $6.1\times10^{-3}$ & $6.1\times10^{-3}$ & $0.010$ & $0.010$ \\
S8 & $7.2\times10^{-6}$ & $1.8\times10^{-14}$ & $7.2\times10^{-3}$ & $7.2\times10^{-3}$ & $0.012$ & $0.012$ \\
H8 & $5.8\times10^{-6}$ & $1.4\times10^{-14}$ & $5.5\times10^{-3}$ & $5.4\times10^{-3}$ & $9.1\times10^{-3}$ & $9.1\times10^{-3}$ \\
\bottomrule
\end{tabular}
\end{table}

\FloatBarrier
\section{P1-HAR-01: data, protocol and cost}
\label{app:har}

\emph{Data.} The official archive (61,005,872 bytes; sha256 recorded with the per-file hashes in the acquisition record) was downloaded once from the UCI repository, checked for path traversal before extraction, and kept outside version control. The repository page states a CC BY 4.0 license; the README requires citing \citet{anguita2013har} and prohibits commercial use; both statements are recorded. Channels: total acceleration (g), body acceleration (g, gravity removed by the distributors' 0.3\,Hz filter) and angular velocity (rad/s). Train 7352 windows / test 2947 windows, as distributed.
\emph{Calibration rule.} The 21 training subjects are ranked by the SHA-256 digest of ``P1-HAR-01:calibration:$\langle$id$\rangle$'' and the first four (7, 15, 17, 19) are the calibration subjects; the runner re-derives the rule.
\emph{Budget rule.} A 20-update smoke run (model discarded; run twice, before and after the pre-execution amendment, 16\,s and $\HarSmokeWall$\,s, both in the ledger) measured the per-update and per-token cost and projected the training and evaluation cost; the pre-registered rule allowed the planned 2000 updates. Ledger: training $\HarTrainWall$\,s wall ($\HarTrainCPU$\,s process CPU), evaluation $\HarEvalWall$\,s wall; all stages together $\HarTotalWall$\,s wall and $\HarTotalCPU$\,s process CPU on 2 threads, against a cap of 3600\,s / 7200\,s.
\emph{Checks.} Eval mode enforced at calibration and test (no dropout, running batch statistics); all step scales positive; no padding or mask; the resampled input equals C0 on S8 and H8 to $\HarResIdent$; the head-of-mean identity $A(\sum_jw_jh_j)=\sum_jw_jA(h_j)$ holds on C0 to $\HarPoolIdent$ and the two native arms coincide on the uniform C0 grid to $\HarArmIdent$; the identities of the logit decomposition hold to $\HarDecIdent$; $\HarExcluded$ windows excluded by the centered-logit rule.
\emph{Timing (\cref{tab:har-timing}).} native\_S8\_end\_to\_end\,/\,resampled\_S8\_end\_to\_end $=\HarRatioSeight$; forward-only $\HarFratioSeight$; H8 end-to-end $\HarRatioHeight$.
\emph{Common-time diagnostic (P1-HAR-CT-01; exploratory, fixed checkpoint, designed after the P1-HAR-01 results).} The stored per-token logits (float16) reproduce the stored pooled logits only to $\HarCtFpSixteen$, so one float32 streaming forward of C0 and S8 over the 2947 test windows was run (batches of 256, eval mode, no intermediate tensors stored); it reproduces the stored C0 and S8 pooled logits to $\HarCtRepro$. The 128 common end points are the S8 tokens $8k$ ($k=1..128$), verified against the time stamps. The split is an algebraic identity of the three logit vectors (the recorded residual, $\HarCtResid$, is zero by construction and checks arithmetic only, not the model); the class centering $\Pi=I-\mathbf{1}\mathbf{1}^{\top}/6$ is applied to every term alike, and the relative norms are per-window ratios averaged subject-first (as a ratio of subject-mean norms the common-time part is $\HarCtNormEnd/\HarCtNormZzero=\HarCtRatioOfMeans$). The pre-specified float64 check (18 hash-selected windows) was \texttt{\HarCtFpSixtyFour} by its rule (threshold $10^{-3}$ on the subject-mean relative common-time change; observed $\HarCtRelEnd$); in absolute terms the common-time change ($\HarCtNormEnd$ logits, subject mean) is about $\HarCtResidualFold$ times the float32 pooling residual recorded by P1-HAR-01. The common-time term is the change of the pooled readout over the 128 end points, not a per-token change: a per-token common-time logit change was not computed (the per-end-point feature change, $\HarCtFeatEnd$ relative, at most $\HarCtFeatEndMax$ in a subject, is recorded but does not establish the per-token logit change). Cost: $\HarCtWall$\,s wall, $\HarCtCPU$\,s process CPU on 2 threads (timer started before importing the framework). Norm sums and norm ratios are not contributions; the inner product is reported alongside.

% AUTO-GENERATED by scripts/make_paper_assets.py from results/raw -- do not edit by hand.
\begin{table*}[t]
\caption{P1-HAR-CT-01 (exploratory, fixed checkpoint, designed after the P1-HAR-01 results): decomposition of the pooled S8 change of the HAR checkpoint, per test subject (subject means first; last row: equal-weight mean). With the official affine head $A$, $z_0=A(\overline{H_0})$ (C0 readout), $z_{\mathrm{end}}=A(\mathrm{mean}_kH_8[\mathrm{end}(k)])$ (S8 features at the 128 common time points), $z_{\mathrm{all}}=A(\overline{H_8})$ (S8 readout), and $z_{\mathrm{all}}-z_0=(z_{\mathrm{end}}-z_0)+(z_{\mathrm{all}}-z_{\mathrm{end}})$ per window exactly. Columns 2--4: norms of the class-centered terms relative to $\|\Pi z_0\|$ (all / common-time / extra-sampling); 5: inner product of the two centered terms (logit$^2$; the norms are not shares); 6: mean relative change of the final features at the common end points; 7--9: cross-entropy of $z_0$, $z_{\mathrm{end}}$, $z_{\mathrm{all}}$; 10--11: fraction of windows whose predicted class differs from $z_0$. No interval and no test: a descriptive diagnostic of one checkpoint.}
\label{tab:har-ct}
\centering\small
\setlength{\tabcolsep}{5pt}
\begin{tabular}{rcccccccccc}
\toprule
subject & rel.\ all & rel.\ common & rel.\ extra & $\langle\cdot,\cdot\rangle$ & feat.\ end & CE $z_0$ & CE $z_{\mathrm{end}}$ & CE $z_{\mathrm{all}}$ & flip end & flip all \\
\midrule
2 & $0.0332$ & $0.0024$ & $0.0321$ & $+0.490$ & $0.0035$ & $0.193$ & $0.193$ & $0.196$ & $0.000$ & $0.003$ \\
4 & $0.0224$ & $0.0015$ & $0.0217$ & $+0.029$ & $0.0015$ & $0.244$ & $0.244$ & $0.280$ & $0.000$ & $0.009$ \\
9 & $0.0313$ & $0.0013$ & $0.0307$ & $+0.125$ & $0.0019$ & $0.455$ & $0.454$ & $0.412$ & $0.000$ & $0.017$ \\
10 & $0.0295$ & $0.0019$ & $0.0288$ & $+0.171$ & $0.0029$ & $1.132$ & $1.133$ & $1.037$ & $0.000$ & $0.031$ \\
12 & $0.0331$ & $0.0020$ & $0.0324$ & $+0.013$ & $0.0027$ & $0.106$ & $0.106$ & $0.092$ & $0.000$ & $0.000$ \\
13 & $0.0309$ & $0.0017$ & $0.0301$ & $+0.100$ & $0.0019$ & $0.119$ & $0.119$ & $0.138$ & $0.000$ & $0.009$ \\
18 & $0.0251$ & $0.0010$ & $0.0247$ & $+0.009$ & $0.0010$ & $0.072$ & $0.072$ & $0.079$ & $0.000$ & $0.003$ \\
20 & $0.0256$ & $0.0036$ & $0.0253$ & $-0.156$ & $0.0054$ & $0.075$ & $0.076$ & $0.074$ & $0.000$ & $0.003$ \\
24 & $0.0303$ & $0.0010$ & $0.0299$ & $+0.038$ & $0.0014$ & $0.037$ & $0.037$ & $0.038$ & $0.000$ & $0.000$ \\
\midrule
mean & $0.0290$ & $0.0018$ & $0.0284$ & $+0.091$ & $0.0025$ & $0.270$ & $0.271$ & $0.261$ & $0.000$ & $0.008$ \\
\bottomrule
\end{tabular}
\end{table*}


\emph{Initialization sensitivity (R8-P1-HAR-REPLICATION).} The run identifier names a sensitivity check on the same subjects and data, not an independent replication. The contract (seeds, stages, inputs with hashes, estimands, budget) was committed before execution. Seeds 101 and 102 were never used before in this project; seed $s$ seeds the parameter initialization (python, NumPy and torch generators) and a separate minibatch-order generator; data, splits and normalization are fixed arrays. Each checkpoint was evaluated once per condition in batches of 256 windows; the S8 common-end-point outputs and the H8 per-token logits were taken from the same forward passes and stored in float64; the resampled arm's input equals C0 exactly (checked), so its logits are the C0 logits and no forward was spent; timing was not re-measured. Selected updates $\RviiiStepA$ and $\RviiiStepB$ (calibration cross-entropy $\RviiiCalCEA$, $\RviiiCalCEB$). Cost: $\RviiiWall$\,s wall and $\RviiiCPU$\,s process CPU for the five stages on 2 threads (caps 1800\,s / 3600\,s); the first training stage ran at $280$\,s wall under concurrent CPU load from bookkeeping commands, within its 300\,s allowance. \Cref{tab:har-r8-risk,tab:har-r8-decomp} give the three runs as separate rows.

% AUTO-GENERATED by scripts/make_paper_assets.py from results/raw -- do not edit by hand.
\begin{table*}[t]
\caption{Initialization sensitivity (run identifier R8-P1-HAR-REPLICATION; a sensitivity check on the same subjects and data, not an independent replication): the frozen P1-HAR-01 protocol trained with two further model seeds (101, 102) on the same training data and evaluated on the SAME nine test subjects; the original run (seed 0) is repeated from \cref{tab:har-subjects,tab:har-effects}. Rows are separate runs and are never pooled (three runs on the same subjects are not 27 independent observations). step / calib.\ CE: selected update and its calibration cross-entropy. CE C0, acc.\ C0: subject-first means of the official readout on C0. $D$: primary difference CE(mean, H8) $-$ CE(time, H8), equal-weight mean over subjects, with the per-run subject-level paired bootstrap interval (2000 draws; coarse with 9 clusters) and the number of subjects with $d_s>0$. $\Delta$CE: cross-entropy change against C0 for S8 (token mean) and H8 (token mean / time-weighted).}
\label{tab:har-r8-risk}
\centering\small
\setlength{\tabcolsep}{4pt}
\begin{tabular}{lccccccccc}
\toprule
seed & step & calib.\ CE & CE C0 & acc.\ C0 & $D$ [95\% CI] & $d_s>0$ & $\Delta$CE S8 & $\Delta$CE H8, mean & $\Delta$CE H8, time \\
\midrule
0 (orig.) & 1200 & $0.092$ & $0.270$ & $0.904$ & $+0.016$ $[-0.004,\,+0.032]$ & 8/9 & $-0.010$ & $+0.012$ & $-0.004$ \\
101 & 1600 & $0.120$ & $0.252$ & $0.916$ & $+0.050$ $[+0.023,\,+0.082]$ & 8/9 & $+0.005$ & $+0.052$ & $+0.002$ \\
102 & 1600 & $0.083$ & $0.252$ & $0.923$ & $+0.044$ $[+0.011,\,+0.077]$ & 7/9 & $-0.026$ & $+0.028$ & $-0.016$ \\
\bottomrule
\end{tabular}
\end{table*}

% AUTO-GENERATED by scripts/make_paper_assets.py from results/raw -- do not edit by hand.
\begin{table*}[t]
\caption{Initialization sensitivity of the decompositions (same three runs as \cref{tab:har-r8-risk}; subject-first means; no interval, no test). S8: class-centered relative change of the pooled logits (``all''), its common-time part $\|\Pi(z_{\mathrm{end}}-z_0)\|/\|\Pi z_0\|$ and extra-sampling part (per-window ratios, subject-first; not additive shares), the inner product of the two centered terms (logit$^2$), and the fraction of windows whose predicted class differs from C0 under $z_{\mathrm{all}}$ and under $z_{\mathrm{end}}$. H8: class-centered relative change of the token-mean and time-weighted readouts, and the norms (logit units) of the token-mean change, its density-weighting term $W$, the model term $M$, its interval-end part $M_{\mathrm{end}}$ and the time-weighted change $Q$.}
\label{tab:har-r8-decomp}
\centering\small
\setlength{\tabcolsep}{4pt}
\begin{tabular}{lccccccccccccc}
\toprule
 & \multicolumn{6}{c}{S8} & \multicolumn{7}{c}{H8} \\
\cmidrule(lr){2-7}\cmidrule(lr){8-14}
seed & all & common & extra & $\langle\cdot,\cdot\rangle$ & flip all & flip end & mean & time & $\|\Delta\|$ & $\|W\|$ & $\|M\|$ & $\|M_{\mathrm{end}}\|$ & $\|Q\|$ \\
\midrule
0 (orig.) & $0.0290$ & $0.0018$ & $0.0284$ & $+0.091$ & $0.008$ & $0.000$ & $0.210$ & $0.016$ & $7.69$ & $7.44$ & $0.997$ & $0.013$ & $0.565$ \\
101 & $0.0231$ & $0.0027$ & $0.0216$ & $+0.215$ & $0.003$ & $0.000$ & $0.242$ & $0.012$ & $9.99$ & $9.60$ & $0.712$ & $0.031$ & $0.457$ \\
102 & $0.0445$ & $0.0034$ & $0.0451$ & $-0.142$ & $0.006$ & $0.000$ & $0.204$ & $0.023$ & $5.92$ & $5.77$ & $1.366$ & $0.016$ & $0.771$ \\
\bottomrule
\end{tabular}
\end{table*}

% AUTO-GENERATED by scripts/make_paper_assets.py from results/raw -- do not edit by hand.
\begin{table}[h]
\caption{P1-HAR-01 timing (one batch of the first 256 test windows; median of 5 timed calls after 2 warm-ups; \texttt{inference\_mode}; 2 threads). End-to-end includes tensor construction, LOCF resampling for the resampled arm, forward and pooling; forward-only is the model call on prebuilt tensors. CPU seconds are process CPU time summed over threads. Timed repeats are not independent samples.}
\label{tab:har-timing}
\centering\small
\begin{tabular}{rrccccc}
\toprule
cond. & tokens & nat.\ e2e & nat.\ fwd & res.\ e2e & res.\ fwd & nat.\ e2e CPU \\
\midrule
C0 & 128 & $0.276$ & $0.306$ & $0.282$ & $0.353$ & $0.55$ \\
S8 & 1024 & $3.562$ & $3.288$ & $0.281$ & $0.266$ & $6.52$ \\
H8 & 576 & $1.359$ & $1.491$ & $0.273$ & $0.273$ & $2.55$ \\
\bottomrule
\end{tabular}
\end{table}


\FloatBarrier
\section{Interval conventions in the TIDES paper and code: details}
\label{app:convention}

The TIDES paper writes $x_{k+1}=\bar\Lambda(\Delta_k)x_k+\bar B(\Delta_k)u_k$ with $\Delta_k=t_{k+1}-t_k$ and $y_k=Cx_k+Du_k$, i.e.\ $u_k$ is held forward on $[t_k,t_{k+1})$ and the state in $y_k$ does not yet contain $u_k$~\citep{soydan2026tides}.
The pinned code scans $z_k=\bar\Lambda(t_k-t_{k-1})z_{k-1}+\bar B(t_k-t_{k-1})u_k$ with $y_k=\mathrm{Re}(Cz_k)+Du_k$, and its forecasting collate builds the step sequence $[t_0,t_1-t_0,\dots]$: $u_k$ is held backward on $(t_{k-1},t_k]$ and the state in $y_k$ contains it.
Both are exact ZOH solutions, but of different input reconstructions.
On a uniform grid every gap is equal, so $z_k=x_{k+1}$: the difference is a one-index shift of the state, which changes the alignment between the state part of $y_k$ and the feedthrough $Du_k$.
On an irregular grid each value is paired with a different gap, so the difference is not a relabeling.
With the pinned code, untrained weights and a fixed impulse input, the code equals the backward-hold recurrence to $\ConvCodeRight$; the one-index shift reproduces the paper form to $\ConvShiftUniform$ on a uniform grid but leaves $\ConvShiftIrregular$ on an irregular one.
The released code, not the printed equations, defines the model behind TIDES' reported results, and it coincides with the right-endpoint convention used throughout this paper; our splitting protocol repeats the right-endpoint value accordingly.
We make no claim that this affects TIDES' results.

\FloatBarrier
\section{Proof of \texorpdfstring{\cref{prop:cascade}}{Proposition 4.2}}
\label{app:proof}

\begin{proof}
(a) follows as in \cref{prop:split}, since the coupled system has constant coefficients on each held interval.
(b) Layer~1 is exact, so with $\tau_k=dt_k$,
$e_k=e^{-c\tau_k}e_{k-1}+d\int_0^{\tau_k}e^{-c(\tau_k-s)}[h_1(t_k)-h_1(t_{k-1}+s)]\,ds$.
If $a>0$ then $|h_1|\le|b|U/a$, so $a|h_1|\le|b|U$; if $a=0$ then $a|h_1|=0$; either way $|\dot h_1|\le2|b|U$.
On each held interval $h_1$ is $C^1$, so the bracket is at most $2|b|U(\tau_k-s)$, and with $e^{-c(\tau_k-s)}\le1$ the local term is at most $|b\,d|U\tau_k^2$.
Unrolling from $e_0=0$ gives $|e_k|\le|b\,d|U\delta\sum_{j\le k}\tau_je^{-c(t_k-t_j)}$, and since $e^{-c(t_k-t_j)}\le e^{c\delta}e^{-c(t_k-s)}$ for $s\in(t_{j-1},t_j]$, the sum is at most $e^{c\delta}\int_0^{t_k}e^{-c(t_k-s)}ds\le e^{c\delta}/c$.
\end{proof}
In the stored P1-COMP-01 outputs, the largest error at the base times is 9--20\% of the bound for $m\in\{1,2,4,8\}$, which is consistent with, not a proof of, the bound.

\FloatBarrier
\section{Numerical floor}
\label{app:numerics}

The naive coefficient $(e^{z}-1)/\lambda$ loses relative precision as $|z|$ shrinks ($\NoneNaive$ at $z{=}-10^{-14}$ in float64), while \texttt{expm1} stays at rounding level (\cref{tab:numerics}).
In emulated float32, split composition of one scalar mode drifts to $\NfourFloor$ at $16{,}384$ sub-steps, so differences below about $10^{-4}$ are not interpretable in that emulation; this is not a floor for a trained multi-layer network with 256 tokens, for which \cref{sec:real} compares float32 with float64 directly.

% AUTO-GENERATED by scripts/make_paper_assets.py from results/raw -- do not edit by hand.
\begin{table}[h]
\caption{Numerical floor (FR-E4). Left: relative error of the ZOH input coefficient $(e^{z}-1)/\lambda$ against a 50-digit reference. Right: splitting $[0,1]$ into $m$ held sub-steps
($\lambda{=}-1$, $g{=}\ln2$, exact $h(1)=0.5$); float32 is an \emph{emulation} (every primitive rounded, correctly rounded $\exp$), not a GPU kernel.}
\label{tab:numerics}
\centering\small
\begin{tabular}{rcc}
\toprule
$z$ & naive & \texttt{expm1} \\
\midrule
$-0.01$ & $5.4\times10^{-15}$ & $4.7\times10^{-17}$ \\
$-1\times10^{-4}$ & $1.3\times10^{-14}$ & $3.4\times10^{-17}$ \\
$-1\times10^{-6}$ & $1.6\times10^{-11}$ & $7.9\times10^{-17}$ \\
$-1\times10^{-8}$ & $1.1\times10^{-9}$ & $6.4\times10^{-17}$ \\
$-1\times10^{-10}$ & $8.3\times10^{-8}$ & $3.4\times10^{-17}$ \\
$-1\times10^{-12}$ & $2.2\times10^{-5}$ & $2.4\times10^{-17}$ \\
$-1\times10^{-14}$ & $8.0\times10^{-4}$ & $4.9\times10^{-17}$ \\
\bottomrule
\end{tabular}\hspace{1em}
\begin{tabular}{rccc}
\toprule
$m$ & fp64 \texttt{expm1} & fp32 \texttt{expm1} & fp32 naive \\
\midrule
$1$ & $0$ & $0$ & $0$ \\
$4$ & $2.2\times10^{-16}$ & $0$ & $6.0\times10^{-8}$ \\
$16$ & $2.2\times10^{-16}$ & $1.2\times10^{-7}$ & $2.4\times10^{-7}$ \\
$64$ & $4.4\times10^{-16}$ & $6.0\times10^{-8}$ & $0$ \\
$256$ & $1.8\times10^{-15}$ & $4.3\times10^{-6}$ & $5.0\times10^{-6}$ \\
$1024$ & $4.4\times10^{-15}$ & $1.8\times10^{-5}$ & $2.0\times10^{-5}$ \\
$4096$ & $1.1\times10^{-13}$ & $2.1\times10^{-5}$ & $2.8\times10^{-5}$ \\
$16384$ & $5.2\times10^{-13}$ & $1.1\times10^{-4}$ & $2.3\times10^{-4}$ \\
\bottomrule
\end{tabular}
\end{table}


\FloatBarrier
\section{Reproducibility}
\label{app:repro}

The anonymized supplementary package contains the code (toy, RK4 references, cascade check, TIDES adapter, convention check, runners, diagnostics and the generator that produces every table, figure and number in this paper), the configurations, and the raw records from which the generator works, together with the sha256 of every checkpoint.
It contains the three 116\,KB UCI HAR checkpoints (model seeds 0, 101 and 102) with their SHA-256 values, but not the synthetic-task checkpoints (their SHA-256 values are listed and they are regenerated by the training command in the package; a retrain of the development seed reproduced every tensor), not the HAR data (acquisition command and archive hash are given) and not the third-party TIDES code, which is obtained at the pinned commit under its own license.

\FloatBarrier
\section{Additional toy tables}
\label{app:toy}

\Cref{tab:interaction,tab:split,tab:coupling,tab:stiff,tab:newobs,tab:readout} list the toy measurements referenced in \cref{sec:toy}: one fix at a time, splitting, the two-layer cascade, the stiff mode, new observations with the exact decomposition, and readouts under a density change.

% AUTO-GENERATED by scripts/make_paper_assets.py from results/raw -- do not edit by hand.
\begin{table}[t]
\caption{One fix at a time (FR-E3-READOUT, split of the first half, $m{=}8$; median scale-normalized readout change over 32 units). The effects are not additive,
so no per-fix share is reported.}
\label{tab:interaction}
\centering\small
\begin{tabular}{lc}
\toprule
fixes applied & change \\
\midrule
none ($\Delta{=}\tau g$, Euler-$B$, sample mean) & $0.46$ \\
time readout only (trapezoidal) & $0.25$ \\
exact ZOH $B$ only & $0.48$ \\
$\Delta{=}dt\,g$ only & $0.30$ \\
$\Delta{=}dt\,g$ + trapezoidal & $0.049$ \\
$\Delta{=}dt\,g$ + exact ZOH (sample mean) & $0.30$ \\
$\Delta{=}dt\,g$ + exact ZOH + trapezoidal & $0.011$ \\
$\Delta{=}dt\,g$ + exact ZOH + exact integral & $1.5\times10^{-16}$ \\
\bottomrule
\end{tabular}
\end{table}

% AUTO-GENERATED by scripts/make_paper_assets.py from results/raw -- do not edit by hand.
\begin{table}[t]
\caption{Splitting a held interval into $m$ sub-steps (FR-E1-SPLIT; toy, 32 units, uniform base grid, real modes).
Entries are the median over units of the output artifact $\|y_m-y^\star\|_2/\|y^\star\|_2$ at the 17 common base times, where $y^\star$ is the exact held-path solution;
slope is the median per-unit log--log slope over $m\in\{1,2,4,8\}$. $\tau$ is the base step, so $\Delta{=}\tau g$ coincides with $\Delta{=}dt\,g$ at $m{=}1$.}
\label{tab:split}
\centering\small
\begin{tabular}{lcccccc}
\toprule
variant & $m{=}1$ & $m{=}2$ & $m{=}4$ & $m{=}8$ & max, $m{=}8$ & slope \\
\midrule
exact ZOH, $\Delta{=}dt\,g$ & $0$ & $2.6\times10^{-16}$ & $3.8\times10^{-16}$ & $5.5\times10^{-16}$ & $6.3\times10^{-15}$ & -- \\
Euler-$B$, $\Delta{=}dt\,g$ & $0.21$ & $0.10$ & $0.050$ & $0.025$ & $0.12$ & $-1.03$ \\
bilinear, $\Delta{=}dt\,g$ & $9.5\times10^{-3}$ & $2.3\times10^{-3}$ & $5.6\times10^{-4}$ & $1.4\times10^{-4}$ & $8.1\times10^{-4}$ & $-2.02$ \\
exact ZOH, $\Delta{=}\tau g$ & $0$ & $0.48$ & $0.93$ & $1.3$ & $2.8$ & -- \\
Euler-$B$, $\Delta{=}\tau g$ & $0.21$ & $0.68$ & $1.1$ & $1.4$ & $3.1$ & $0.93$ \\
bilinear, $\Delta{=}\tau g$ & $9.5\times10^{-3}$ & $0.48$ & $0.94$ & $1.3$ & $2.8$ & $2.25$ \\
\bottomrule
\end{tabular}
\end{table}

% AUTO-GENERATED by scripts/make_paper_assets.py from results/raw -- do not edit by hand.
\begin{table}[t]
\caption{Two-layer linear cascade $\dot h_1=-ah_1+bu$, $\dot h_2=-ch_2+dh_1$ under interval splitting (P1-COMP-01; $b{=}d{=}1$; held input of FIRST\_RUN unit 0).
Columns $m$: error of the layerwise computation (layer 2 holds the layer-1 endpoint value) against the exact coupled solution, $\|h_2^{\mathrm{lw}}-h_2\|/\|h_2\|$ over base times.
``change'' is the layerwise change $m{=}8$ vs.\ $m{=}1$; the last column is the largest change of the exact coupled solution over $m$. No external observation is added in this experiment.}
\label{tab:coupling}
\centering\small
\begin{tabular}{lccccccc}
\toprule
case & $m{=}1$ & $m{=}2$ & $m{=}4$ & $m{=}8$ & slope & change & exact inv. \\
\midrule
$a{=}1,\;c{=}0.5$ & $0.19$ & $0.10$ & $0.052$ & $0.026$ & $-0.95$ & $0.15$ & $1.7\times10^{-15}$ \\
$a{=}c{=}1$ & $0.20$ & $0.11$ & $0.055$ & $0.028$ & $-0.93$ & $0.16$ & $7.1\times10^{-16}$ \\
\bottomrule
\end{tabular}
\end{table}

% AUTO-GENERATED by scripts/make_paper_assets.py from results/raw -- do not edit by hand.
\begin{table}[t]
\caption{Stiff scalar mode (FR-E4-N2; $g{=}\ln 2$, base step $0.5$, held sine input). Bilinear is stable ($|\bar A|<1$) for every step, but for $z{=}\Delta\lambda<-2$ its transition is negative
(sign-alternating transients) and tends to $-1$, whereas the exact decay $e^{z}$ tends to $0$. The last two columns are the output artifact against exact ZOH.
The stable/sign/decay columns are an exploratory re-reading of logged values.}
\label{tab:stiff}
\centering\small
\begin{tabular}{rrrrrccrr}
\toprule
$\lambda$ & $m$ & $z$ & $\bar A_{\mathrm{bil}}$ & $e^{z}$ & stable & sign-alt. & bilinear & Euler-$B$ \\
\midrule
$-5$ & $1$ & $-1.73$ & $+0.072$ & $0.18$ & yes & no & $0.098$ & $1.1$ \\
$-5$ & $4$ & $-0.43$ & $+0.644$ & $0.65$ & yes & no & $4.7\times10^{-3}$ & $0.23$ \\
$-5$ & $16$ & $-0.11$ & $+0.897$ & $0.90$ & yes & no & $2.9\times10^{-4}$ & $0.055$ \\
$-5$ & $32$ & $-0.05$ & $+0.947$ & $0.95$ & yes & no & $7.3\times10^{-5}$ & $0.027$ \\
$-50$ & $1$ & $-17.33$ & $-0.793$ & $3.0\times10^{-8}$ & yes & yes & $0.54$ & $16$ \\
$-50$ & $4$ & $-4.33$ & $-0.368$ & $0.013$ & yes & yes & $0.015$ & $3.4$ \\
$-50$ & $16$ & $-1.08$ & $+0.297$ & $0.34$ & yes & no & $2.1\times10^{-8}$ & $0.64$ \\
$-50$ & $32$ & $-0.54$ & $+0.574$ & $0.58$ & yes & no & $8.5\times10^{-9}$ & $0.30$ \\
$-500$ & $1$ & $-173.29$ & $-0.977$ & $5.5\times10^{-76}$ & yes & yes & $1.1$ & $170$ \\
$-500$ & $4$ & $-43.32$ & $-0.912$ & $1.5\times10^{-19}$ & yes & yes & $0.62$ & $42$ \\
$-500$ & $16$ & $-10.83$ & $-0.688$ & $2.0\times10^{-5}$ & yes & yes & $2.0\times10^{-3}$ & $9.8$ \\
$-500$ & $32$ & $-5.42$ & $-0.461$ & $4.4\times10^{-3}$ & yes & yes & $1.3\times10^{-11}$ & $4.4$ \\
\bottomrule
\end{tabular}
\vspace{4pt}

\begin{tabular}{rccc}
\toprule
$dt$ & exact ZOH & Euler-$B$ & bilinear \\
\midrule
$0.01$ & $8.0\times10^{-15}$ & $3.5\times10^{-3}$ & $0$ \\
$0.5$ & $0$ & $0.18$ & $1.1\times10^{-16}$ \\
$8$ & $0$ & $4.6$ & $5.6\times10^{-6}$ \\
\bottomrule
\end{tabular}
\vspace{2pt}

{\footnotesize Lower block (FR-E4-N3): relative steady-state error under a held unit input, $\lambda{=}-1$.}
\end{table}

% AUTO-GENERATED by scripts/make_paper_assets.py from results/raw -- do not edit by hand.
\begin{table}[t]
\caption{Adding new observations of the same signal (FR-E2-NEWOBS; toy, 32 units). The change $T=y_m-y_1$ at base times splits \emph{exactly} as $T=P+R$, where $P$ is the change of the continuous-time solution on the new held path
(the change in the held-input reconstruction; no task or Bayesian information gain is implied) and $R$ is the change of the discretization artifact. Hence $\|T\|^2=\|P\|^2+\|R\|^2+2\langle P,R\rangle$; we report these three terms as fractions of $\sum_{\text{units}}\|T\|^2$ (they sum to 1)
and the per-unit range of the cross fraction. $P$ uses the exact-ZOH output as a proxy for the continuous-time held-path solution (validated against RK4 to the bound shown).
Median total change is relative to $\|y_1\|$. EXPLORATORY re-aggregation; the cross term shows that no per-unit ``share'' is well defined, and none of these fractions is a causal contribution.}
\label{tab:newobs}
\centering\small
\begin{tabular}{lcccccc}
\toprule
variant & $m$ & median $\|T\|/\|y_1\|$ & $\|P\|^2$ & $\|R\|^2$ & $2\langle P,R\rangle$ & cross range \\
\midrule
exact ZOH, $\Delta{=}dt\,g$ & $2$ & $0.11$ & \multicolumn{4}{c}{artifact $\le1.7\times10^{-13}$ vs.\ RK4 (proxy reference)} \\
exact ZOH, $\Delta{=}dt\,g$ & $8$ & $0.19$ & \multicolumn{4}{c}{artifact $\le1.7\times10^{-13}$ vs.\ RK4 (proxy reference)} \\
Euler-$B$, $\Delta{=}dt\,g$ & $2$ & $0.15$ & $0.173$ & $0.567$ & $+0.260$ & $[-0.93,+0.42]$ \\
Euler-$B$, $\Delta{=}dt\,g$ & $8$ & $0.25$ & $0.198$ & $0.533$ & $+0.269$ & $[-1.37,+0.42]$ \\
bilinear, $\Delta{=}dt\,g$ & $2$ & $0.11$ & $0.849$ & $0.025$ & $+0.126$ & $[-1.04,+0.34]$ \\
bilinear, $\Delta{=}dt\,g$ & $8$ & $0.19$ & $0.886$ & $0.013$ & $+0.101$ & $[-0.69,+0.29]$ \\
exact ZOH, $\Delta{=}\tau g$ & $2$ & $0.42$ & $0.071$ & $1.260$ & $-0.331$ & $[-1.06,+0.34]$ \\
exact ZOH, $\Delta{=}\tau g$ & $8$ & $1.1$ & $0.025$ & $1.223$ & $-0.248$ & $[-1.02,+0.03]$ \\
Euler-$B$, $\Delta{=}\tau g$ & $2$ & $0.42$ & $0.053$ & $1.225$ & $-0.278$ & $[-0.72,+0.32]$ \\
Euler-$B$, $\Delta{=}\tau g$ & $8$ & $1.1$ & $0.020$ & $1.202$ & $-0.222$ & $[-0.75,+0.04]$ \\
bilinear, $\Delta{=}\tau g$ & $2$ & $0.42$ & $0.072$ & $1.260$ & $-0.331$ & $[-1.10,+0.35]$ \\
bilinear, $\Delta{=}\tau g$ & $8$ & $1.1$ & $0.025$ & $1.222$ & $-0.247$ & $[-1.04,+0.03]$ \\
\bottomrule
\end{tabular}
\end{table}

% AUTO-GENERATED by scripts/make_paper_assets.py from results/raw -- do not edit by hand.
\begin{table}[t]
\caption{Readouts under a density change (FR-E3-READOUT; toy, 32 units). Columns 2--3: scale-normalized change $|r_8-r_1|/s$ when only the first half of the base intervals is split into $8$ held sub-steps
(physical path unchanged); $s$ is the RMS base-time output at $m{=}1$ ($\times K_1$ for the sample sum). Column 4: scale-normalized error to the continuous-time truth when new observations are added to the first half only, $m{=}1\to8$.
Medians over units. The exact held-path integral is defined only for exact-ZOH states.}
\label{tab:readout}
\centering\small
\begin{tabular}{lccc}
\toprule
readout & split, exact ZOH $\Delta{=}dt\,g$ & split, Euler-$B$ $\Delta{=}\tau g$ & new obs., exact ZOH \\
\midrule
last sample & $8.9\times10^{-17}$ & $0.19$ & $0.085$$\to$$0.099$ \\
sample sum & $1.5$ & $2.9$ & -- \\
sample mean & $0.30$ & $0.46$ & $0.060$$\to$$0.28$ \\
right-endpoint $\sum y_k\,dt_k$ & $0.026$ & $0.24$ & $0.060$$\to$$0.043$ \\
trapezoidal endpoint sum & $0.011$ & $0.25$ & $0.039$$\to$$0.023$ \\
exact held-path integral & $1.5\times10^{-16}$ & n/a & $0.039$$\to$$0.032$ \\
\bottomrule
\end{tabular}
\end{table}


\FloatBarrier
\end{document}
